\chapter{ORGANIZAÇÃO DIDÁTICO-PEDAGÓGICA DO CURSO}

A organização didático-pedagógica do curso está fundamentada no paradigma contemporâneo da formação por competências, entendida como a capacidade de mobilizar, integrar e aplicar saberes teóricos, práticos, socioemocionais e contextuais para resolver problemas reais em ambientes complexos e em transformação. Em consonância com os Referenciais de Formação da Sociedade Brasileira de Computação (SBC) \cite{SBC_2017}, com a Resolução CNE/CES n$^\circ$ 5/2016 \cite{MEC_2016} que define as diretrizes curriculares para cursos na área de Computação, com o PPI \cite{UTFPR_PPI_2019} e com o PDI 2023-2027 da UTFPR \cite{UTFPR_PDI_2023}, o curso orienta-se para a formação de profissionais com autonomia intelectual, responsabilidade ética, criatividade, pensamento crítico e domínio operacional das tecnologias digitais e inteligentes.

A materialização desta proposta ocorre por meio de um conjunto articulado de estratégias acadêmicas e institucionais, que incluem a inserção dos estudantes, desde os períodos iniciais em atividades de pesquisa, iniciação científica e inovação; o desenvolvimento de projetos integradores e de ações de resolução de problemas reais, em diálogo com o setor produtivo e com o ecossistema regional de inovação; a participação ativa em ações de extensão universitária, assegurando a indissociabilidade entre ensino, pesquisa e extensão; e o incentivo ao empreendedorismo tecnológico e à criação de soluções inovadoras, por meio de editais institucionais, programas de pré-incubação, hackathons, desafios tecnológicos e projetos desenvolvidos em parceria com empresas, startups e órgãos públicos.

Do ponto de vista da governança acadêmica, a implementação e atualização contínua do curso são assegurados pela atuação integrada do Núcleo Docente Estruturante (NDE) e do Colegiado de Curso, responsáveis por garantir a coerência curricular, o alinhamento às diretrizes institucionais e nacionais, a atualização periódica do PPC e a aderência às demandas acadêmicas e profissionais emergentes. O estágio curricular supervisionado e o Trabalho de Conclusão de Curso (TCC) constituem momentos privilegiados de síntese e aplicação das competências desenvolvidas ao longo do percurso formativo, possibilitando ao estudante integrar teoria, prática e ação profissional.

Apoiado em políticas institucionais de pesquisa, iniciação científica, extensão, inovação e empreendedorismo, e materializado por meio de práticas pedagógicas intencionalmente planejadas, o curso de CDIA no campus Londrina apresenta uma proposta formativa coerente, contemporânea e socialmente comprometida. Trata-se de uma formação orientada não apenas para acompanhar as transformações tecnológicas em curso, mas para liderá-las, com responsabilidade, visão sistêmica, rigor técnico e compromisso com o desenvolvimento sustentável e com a transformação digital da sociedade.

\section{POLÍTICA INSTITUCIONAL NO ÂMBITO DO CURSO}

O curso de CDIA no campus Londrina foi concebido a partir de estudos institucionais sistemáticos, diagnósticos regionais e análises de demandas nacionais associadas aos setores de tecnologia, inteligência artificial, Indústria 4.0, serviços digitais, saúde conectada, segurança cibernética e desenvolvimento de sistemas inteligentes, conforme detalhado na seção Integração do currículo com o mundo do trabalho (veja Seção 2.1.1). Sua oferta está diretamente vinculada ao papel estratégico da UTFPR como instituição de referência na neo-industrialização brasileira, reconhecimento consolidado no âmbito do programa Universidades Inovadoras do Ministério da Educação.

Em consonância com o Plano de Desenvolvimento Institucional (PDI 2023-2027) \cite{UTFPR_PDI_2023}, o curso materializa a missão da UTFPR de desenvolver educação tecnológica de excelência e produzir conhecimento orientado à solução de desafios reais da sociedade. Nesse contexto, o curso incorpora e operacionaliza as políticas institucionais de ensino, pesquisa, extensão, inovação, empreendedorismo e internacionalização, articulando-as de forma integrada às diretrizes nacionais e às demandas loco-regionais.

A UTFPR apresenta ampla e consolidada inserção em ecossistemas de inovação, incluindo parques tecnológicos, incubadoras, arranjos produtivos locais (APLs), polos de inovação e setores industriais estratégicos do estado do Paraná. Essa articulação institucional possibilita aos estudantes vivenciar situações autênticas de aprendizagem em contextos reais e complexos. O curso de CDIA dialoga intensamente com esses ambientes por meio de práticas profissionais supervisionadas, projetos tecnológicos integradores, pesquisa aplicada, convênios e parcerias institucionais, bem como acordos de cooperação técnica, os quais serão detalhados nas seções correspondentes deste Projeto Pedagógico.

No âmbito do ensino, o curso implementa as políticas institucionais da UTFPR assegurando coerência com o PDI \cite{UTFPR_PDI_2023} e com o Projeto Pedagógico Institucional (PPI) \cite{UTFPR_PPI_2019}. A organização didático-pedagógica adota o modelo UTFPR (STEM$^{IA}$), que articula Ciência, Tecnologia, Engenharia e Matemática ao uso estratégico, crítico e ético da Inteligência Artificial, estruturando percursos formativos orientados por competências e centrados na aprendizagem ativa. A articulação entre teoria e prática é promovida por meio de unidade curriculares integradoras, práticas laboratoriais, projetos baseados em problemas reais, práticas profissionais supervisionadas e participação dos estudantes em eventos técnicos e atividades de vivência profissional.

A implementação dessas políticas é assegurada por condições institucionais adequadas, incluindo a atuação docente em Núcleo Docente Estruturante (NDE), Colegiado de Curso, comissões acadêmicas, orientação de estágios e Trabalho de Conclusão de Curso. O curso conta, ainda, com políticas institucionais de apoio à docência, incluindo ações de acolhimento, acompanhamento e formação continuada de docentes iniciantes, atualização pedagógica, desenvolvimento profissional e incentivo à inovação educacional.

O curso incorpora as políticas institucionais de pesquisa e iniciação científica, estimulando a participação de estudantes em atividades investigativas desde os primeiros períodos. Essa integração se traduz em envolvimento em projetos de pesquisa aplicada e tecnológica, participação em programas de iniciação científica (PIBIC, PIBITI, BIC/UTFPR), desenvolvimento de experimentos, protótipos e soluções inovadoras vinculadas à área do curso e estímulo à publicação científica e à participação em eventos acadêmicos e tecnológicos. Essas ações fortalecem o desenvolvimento de competências investigativas e reforçam a identidade tecnológica e inovadora do curso.

Em consonância com a Política Nacional de Extensão, o curso integra práticas de extensão universitária como componentes essenciais da formação, atendendo à carga horária regulamentar e promovendo vínculo permanente entre universidade e sociedade. As atividades de extensão ocorrem sob a forma de programas, projetos, cursos, oficinas, eventos e prestação de serviços tecnológicos, devidamente registrados nos sistemas institucionais, garantindo rastreabilidade, avaliação e comprovação da carga horária extensionista, além de promover a participação dos estudantes em ações vinculadas a problemas reais e demandas sociais.

O conjunto de políticas de ensino, pesquisa, iniciação científica e extensão implantadas no curso dialoga diretamente com demandas socioeconômicas e tecnológicas do Paraná, especialmente nos setores de automação, manufatura avançada, agritech, saúde digital, segurança cibernética, cidades inteligentes e energias renováveis. Tais políticas também estão enquadradas no contexto da Região Metropolitana de Londrina, que é um importante polo nas áreas de serviços, comércio atividades da cadeia do agronegócio, eletro-metal-mecânica, químicos e materiais, saúde e comunicação. Tais setores fazem uso de tecnologias de informação para inovar em seus produtos e serviços, necessitando de mão de obra altamente especializada e que domine as mais recentes tecnologias.

Essa articulação do Projeto Pedagógico do Curso com as demandas regionais contribui para a inovação em processos produtivos locais, o desenvolvimento sustentável, a redução de desigualdades tecnológicas e o fortalecimento de arranjos produtivos locais e ecossistemas de inovação, evidenciando o impacto social e regional da formação ofertada.

\subsection{Integração do curso com o mundo do trabalho}
\label{sec:integracaoMundoDoTrabalho}

A região de Londrina, localizada no Norte do Paraná, apresenta relevância estratégica no contexto estadual e nacional, consolidada como um dos principais polos econômicos, educacionais e tecnológicos do interior do Brasil \cite{Pinto_2022}. O município exerce influência direta sobre um amplo conjunto de cidades da região norte do estado, configurando-se como um centro regional de serviços avançados, inovação, formação de recursos humanos e articulação produtiva. Trata-se de um território marcado por elevada diversificação econômica, com forte presença dos setores de serviços, comércio, indústria, agronegócio tecnológico e Tecnologia da Informação e Comunicação (TIC) \cite{ibgeCensoLondrina2022}.

Nesse contexto, Londrina consolidou-se como um dos principais polos regionais de TIC do estado do Paraná, destacando-se pela densidade de empresas de base tecnológica e pela articulação entre setor produtivo, instituições de ensino e pesquisa, ambientes de inovação e poder público \cite{Abes_2021}. O setor de TIC exerce papel estruturante na economia local, ao apoiar processos de transformação digital e inovação em diversos segmentos produtivos.

Elemento central desse ecossistema é o Arranjo Produtivo Local de Tecnologia da Informação e Comunicação (APL TIC Londrina), fundado em 2006 e reconhecido como o primeiro APL de TIC do Paraná e um dos pioneiros no Brasil. O APL TIC Londrina tem como missão promover o desenvolvimento e a competitividade do setor de TIC regional por meio da integração entre empresas, instituições de ensino e pesquisa, órgãos governamentais e demais atores da cadeia produtiva.

Complementarmente, o ecossistema de inovação de Londrina é fortalecido pela presença de ambientes estruturantes como a Estação 43, hub de inovação que articula startups, empresas, instituições de ensino, poder público e organizações de apoio ao empreendedorismo. Além disso, o município conta com um Masterplan 2040, que orienta as estratégias de desenvolvimento urbano, econômico e tecnológico, reconhecendo Londrina como polo regional em setores que utilizam intensivamente tecnologias de informação, análise de dados e soluções digitais avançadas.

A crescente adoção de Ciência de Dados e Inteligência Artificial no âmbito do setor de TIC, bem como nos segmentos industrial, agroindustrial, financeiro, logístico, educacional e de serviços, evidencia uma demanda contínua por profissionais altamente qualificados \cite{Brasscom_2021}. Tais profissionais são requeridos para atuar em atividades como análise e engenharia de dados, aprendizado de máquina, sistemas inteligentes, automação de processos analíticos, inteligência de negócios, governança de dados e suporte à tomada de decisão estratégica.

Nesse cenário, a implantação do curso de CDIA no campus Londrina da UTFPR representa uma vantagem estratégica para o município e sua região de influência. O curso alinha a formação acadêmica às demandas concretas do ecossistema produtivo local, em especial do APL TIC Londrina e dos demais setores econômicos estratégicos da região.

A proposta formativa favorece a inserção regional e no mundo do trabalho por meio de projetos integradores, estágios curriculares, atividades de extensão, iniciação científica e tecnológica e parcerias com empresas do APL TIC, startups, incubadoras e ambientes de inovação como a Estação 43. Essas ações promovem a transferência de conhecimento e tecnologia para a sociedade, permitindo que os estudantes vivenciar problemas reais do contexto regional ao longo de sua formação.

Ao formar profissionais aptos a desenvolver, implementar e gerenciar soluções baseadas em Ciência de Dados e Inteligência Artificial, o curso contribui diretamente para o fortalecimento do ecossistema de inovação de Londrina, para o aumento da competitividade regional e para o desenvolvimento socioeconômico sustentável da região Norte do Paraná.

\subsection{Articulação entre Teoria e Prática}

A articulação entre teoria e prática constitui princípio estruturante do curso de CDIA do campus Londrina. Embora concebido a partir de modelos tradicionais de organização curricular, o curso incorpora progressivamente a formação por competências, promovendo a integração entre fundamentos conceituais, desenvolvimento técnico e aplicação prática em contextos reais.

As unidades curriculares Oficina de Integração 1 e Oficina de Integração 2 desempenham papel central nesse processo. Estruturadas com foco em projetos extensionistas e resolução de problemas reais, essas disciplinas possibilitam ao estudante aplicar, de forma sistêmica e interdisciplinar, os conhecimentos adquiridos ao longo do curso. A Oficina de Integração 1 concentra-se na consolidação e integração de conteúdos da área de Computação, enquanto a Oficina de Integração 2 articula conhecimentos específicos de Ciência de Dados e Inteligência Artificial.

Além disso, as unidades curriculares com carga horária prática promovem a vivência de situações que simulam ou reproduzem desafios do mundo do trabalho. Nessas atividades, o docente pode adotar metodologias ativas, estudos de caso, aprendizagem baseada em projetos e outras estratégias didáticas que incentivam o protagonismo discente e a construção significativa do conhecimento.

O Trabalho de Conclusão de Curso (TCC) também se configura como espaço privilegiado de integração entre teoria e prática, ao exigir o desenvolvimento de projetos que mobilizam múltiplas áreas do conhecimento \cite{UTFPR_TCC_2018}.

Por fim, o Estágio Curricular Supervisionado consolida essa articulação ao inserir o estudante em ambientes profissionais reais, nos quais os conhecimentos teóricos são confrontados com situações concretas, promovendo a integração de saberes e o desenvolvimento de competências técnicas \cite{LDE_2008, Estagio_2020}.

Dessa forma, o curso assegura que a formação acadêmica esteja permanentemente conectada à prática profissional, fortalecendo a preparação do egresso para atuar de maneira crítica, competente e inovadora.

\subsection{Articulação com a Pesquisa e Pós-Graduação}

A UTFPR compreende a pesquisa, a iniciação científica e a inovação tecnológica, artística e cultural como um conjunto indissociável de ações voltadas à produção de novos conhecimentos e ao desenvolvimento científico, tecnológico e social. Essas atividades constituem pilares fundamentais da formação acadêmica, na medida em que a prática investigativa promove o pensamento crítico, a reflexão sistemática e a superação da reprodução acrítica de saberes e práticas tradicionais, conforme destacado por Pizzato et al. (2009).

A indissociabilidade entre ensino e pesquisa favorece a construção de uma relação transformadora entre a universidade e a sociedade, possibilitando que estudantes, em diferentes estágios formativos, participem de projetos de investigação apoiados por grupos de pesquisa e pela infraestrutura institucional. Essa articulação fortalece tanto a formação acadêmica quanto a produção de conhecimento orientada às demandas sociais e ao mundo do trabalho, respeitando o rigor metodológico e científico característico da universidade pública.

Nesse contexto, o curso de CDIA do campus Londrina tem como um de seus objetivos estruturantes a consolidação da produção de conhecimento científico e tecnológico, ampliando sua maturidade acadêmica e contribuindo para o desenvolvimento regional e nacional. O curso promove, de forma sistemática, a articulação entre a graduação e a pós-graduação, em consonância com as políticas institucionais de pesquisa da UTFPR e com sua missão de ofertar educação tecnológica de excelência por meio do ensino, da pesquisa e da extensão.

A UTFPR mantém programas institucionais consolidados de Iniciação Científica e Iniciação em Desenvolvimento Tecnológico e Inovação, além de ações afirmativas voltadas à inclusão social. Soma-se a isso a atuação de docentes do curso em Programas de Pós-Graduação stricto sensu da própria instituição e de outras instituições de ensino e pesquisa, o que amplia as possibilidades de inserção dos estudantes em ambientes qualificados de investigação científica e tecnológica, fortalecendo a integração entre ensino, pesquisa e inovação.

Os estudantes do curso são incentivados a participar ativamente de projetos de pesquisa desenvolvidos por docentes do campus Londrina, tanto por meio da integralização de atividades complementares quanto pelo envolvimento direto em projetos vinculados às áreas básicas e específicas da formação em Ciência de Dados e Inteligência Artificial. Essa participação favorece a aplicação prática dos conhecimentos adquiridos, o desenvolvimento do pensamento científico, a autonomia intelectual e a aproximação com problemas reais de natureza acadêmica, tecnológica e social.

Destacam-se, nesse cenário, o Programa Institucional de Bolsas de Iniciação Científica (PIBIC) e o Programa Institucional de Bolsas de Iniciação em Desenvolvimento Tecnológico e Inovação (PIBITI), bem como suas modalidades voluntárias, apoiados por agências de fomento como o CNPq, a Fundação Araucária e a própria UTFPR. Os estudantes podem participar desses programas na condição de bolsistas ou voluntários, sempre vinculados a projetos e grupos de pesquisa formalmente constituídos e certificados.

A atuação em pesquisa está diretamente alinhada ao perfil do egresso do curso, ao contribuir para o desenvolvimento de competências como a identificação e análise de problemas complexos, a busca e sistematização de informações, a formulação de hipóteses, a seleção de métodos adequados, a coleta e análise de dados, a interpretação de resultados e a comunicação de conclusões fundamentadas em evidências científicas. Essas experiências favorecem a formação de profissionais críticos, investigativos e aptos a lidar com situações novas e complexas, além de estimular a continuidade da formação acadêmica em nível de pós-graduação.

As atividades de iniciação científica e tecnológica desenvolvidas no âmbito do curso ocorrem de forma articulada a projetos e grupos de pesquisa certificados, assegurando acompanhamento docente qualificado e rigor metodológico. Como espaço institucional de socialização, avaliação e divulgação dessas produções, a UTFPR promove anualmente o Seminário de Iniciação Científica e Tecnológica (SICITE), organizado pela Pró-Reitoria de Pesquisa e Pós-Graduação (PROPPG), em conjunto com as Diretorias de Pesquisa e Pós-Graduação (DIRPPG) e o Comitê Interno PIBIC/PIBITI.

O SICITE constitui um ambiente acadêmico fundamental para a apresentação e divulgação dos trabalhos desenvolvidos nos programas institucionais de iniciação científica e tecnológica, incluindo PIBIC, PIBIC-Ações Afirmativas, PIBITI, PIBIC Júnior, PIBIC-EM, PIVIC e PIVIT, possibilitando a interação dos estudantes com a comunidade acadêmica, o aprimoramento de suas produções e a difusão do conhecimento gerado à sociedade.

Dessa forma, ao integrar de maneira efetiva as políticas institucionais de pesquisa e iniciação científica à sua proposta formativa, o curso de CDIA assegura uma formação tecnológica de base científica sólida, orientada ao desenvolvimento de competências técnicas, investigativas e inovadoras demandadas por setores produtivos intensivos em tecnologia, reafirmando o compromisso da UTFPR com o avanço do conhecimento, o desenvolvimento social e a transformação da realidade por meio da educação pública de qualidade.

\subsection{Articulação com Políticas de Extensão}

Em consonância com as políticas institucionais da UTFPR, a articulação com as políticas de extensão no curso de CDIA está fundamentada na indissociabilidade entre ensino, pesquisa, iniciação científica e inovação tecnológica, artística e cultural, compreendidas como pilares da atividade acadêmica e como meios privilegiados para a produção de novos conhecimentos. A prática extensionista, nesse contexto, dialoga diretamente com a pesquisa ao incentivar a superação da reprodução acrítica de práticas tradicionais, promovendo processos reflexivos, interação social qualificada e o desenvolvimento de competências metacognitivas, essenciais à formação acadêmica e profissional.

A integração entre ensino, pesquisa e extensão viabiliza uma relação transformadora entre a universidade e a sociedade, na medida em que projetos extensionistas são concebidos e desenvolvidos em articulação com grupos de pesquisa e com o uso compartilhado da infraestrutura institucional, acolhendo estudantes em diferentes estágios formativos. Essa abordagem favorece a construção de soluções fundamentadas no rigor acadêmico e orientadas para o enfrentamento de problemas reais, contribuindo simultaneamente para a formação do estudante e para o atendimento às demandas sociais e do mundo do trabalho.

No âmbito do curso, a extensão universitária constitui espaço privilegiado para a aplicação prática do conhecimento científico e tecnológico produzido nos projetos de pesquisa, permitindo que os estudantes articulem teoria e prática em contextos concretos. As atividades extensionistas favorecem, assim, a dupla conscientização destacada pelas políticas institucionais: de um lado, a do pesquisador em formação, que incorpora como desafio acadêmico a busca por soluções socialmente relevantes; de outro, a da sociedade e dos setores produtivos, que se beneficiam de conhecimentos desenvolvidos sob critérios de qualidade, ética e rigor científico.

Os docentes do curso são incentivados a estruturar e homologar projetos institucionais que integrem pesquisa e extensão, possibilitando a participação de estudantes de graduação como bolsistas ou voluntários de iniciação científica e tecnológica, mediante planos de atividades que contemplem investigação, desenvolvimento e aplicação do conhecimento. Esses projetos frequentemente envolvem revisão da literatura, coleta e análise de dados, desenvolvimento de sistemas e soluções baseadas em Ciência de Dados e Inteligência Artificial, além da produção de artigos, relatórios técnicos e outros produtos acadêmicos e tecnológicos.

As experiências resultantes dessas ações integradas são socializadas em espaços institucionais de divulgação científica e tecnológica, como o Seminário de Iniciação Científica e Tecnológica (SICITE) da UTFPR, promovido anualmente pela Pró-Reitoria de Pesquisa e Pós-Graduação (PROPPG), pelas Diretorias de Pesquisa e Pós-Graduação dos campi e pelo Comitê Interno PIBIC/PIBITI. A participação dos estudantes nesse evento amplia a visibilidade das ações extensionistas articuladas à pesquisa, fortalece a formação acadêmica e contribui para a difusão do conhecimento produzido para a comunidade acadêmica e para a sociedade.

Dessa forma, a articulação com as políticas de extensão no curso de CDIA reafirma o compromisso institucional da UTFPR com uma formação tecnológica de excelência, na qual a extensão se configura como espaço de aplicação social do conhecimento científico, de integração com a pesquisa e de fortalecimento da relação entre universidade, sociedade e mundo do trabalho.

\subsection{Articulação com as Políticas Institucionais de Inovação, Empreendedorismo, Desenvolvimento Tecnológico, Ambiental e Social}

A articulação do curso de CDIA com as políticas institucionais de inovação, empreendedorismo e desenvolvimento tecnológico, ambiental e social da Universidade Tecnológica Federal do Paraná (UTFPR) fundamenta-se na compreensão de que a formação tecnológica contemporânea deve integrar, de maneira indissociável, competência técnica, atitude empreendedora, responsabilidade socioambiental e compromisso ético com o desenvolvimento sustentável. Nesse sentido, o curso adota uma abordagem formativa orientada à inovação responsável, à aplicação socialmente relevante do conhecimento e à construção de soluções tecnológicas alinhadas às demandas da sociedade e do mundo do trabalho.

Em consonância com o Plano de Desenvolvimento Institucional \cite{UTFPR_PDI_2023}, com o Projeto Pedagógico Institucional \cite{UTFPR_PPI_2019}, com as Diretrizes Curriculares dos Cursos de Graduação Regulares da UTFPR \cite{UTFPR_DC_2022} e com a Política de Sustentabilidade da UTFPR, aprovada pela Deliberação COUNI n$^o$ 07/2019 \cite{UTFPR_SUST_POL_2019}, o curso promove a inovação como prática formativa transversal. Essa diretriz orienta a integração entre conhecimentos científicos e tecnológicos e a análise crítica de seus impactos ambientais, sociais e econômicos, contribuindo para a formação de profissionais capazes de atuar de forma ética, crítica e socialmente responsável. Conforme o art. 3$^o$ da Resolução COGEP n$^o$ 142/2022 \cite{UTFPR_DC_2022}, o Projeto Pedagógico do Curso enfatiza a vivência dos estudantes com problemas reais da sociedade, especialmente aqueles relacionados ao desenvolvimento socioeconômico, à aplicação da tecnologia e à busca de alternativas inovadoras para a resolução de desafios técnicos, organizacionais e sociais.

As políticas institucionais de inovação e empreendedorismo da UTFPR são operacionalizadas por meio da integração do curso com programas, estruturas e ações coordenadas pela Pró-Reitoria de Relações Empresariais e Comunitárias (PROREC), pela Diretoria de Relações Empresariais e Comunitárias (DIREC) e pelo Núcleo de Inovação Tecnológica (NIT), com destaque para o Programa de Empreendedorismo e Inovação da UTFPR (PROEM). Essas instâncias fomentam a cultura empreendedora, o desenvolvimento de projetos inovadores, a proteção da propriedade intelectual, a transferência de tecnologia e a aproximação entre universidade, setor produtivo e sociedade, incluindo iniciativas voltadas à inovação responsável, à economia criativa, à economia circular e à geração de impacto social positivo.

No âmbito do campus Londrina, o curso articula-se diretamente com o ecossistema local de inovação, com o Arranjo Produtivo Local (APL) de Tecnologia da Informação e Comunicação (TIC) de Londrina e com a Incubadora Sprint do campus Londrina da UTFPR. Essa articulação possibilita o desenvolvimento de soluções tecnológicas baseadas em Ciência de Dados e Inteligência Artificial com impacto social, ambiental e econômico, fortalecendo a inserção regional do curso e ampliando as oportunidades formativas dos estudantes em ambientes reais de inovação, empreendedorismo e desenvolvimento tecnológico.

Destaca-se, ainda, a existência no campus do Laboratório Multiusuário do campus Londrina (LabMult-LD), das empresas juniores vinculadas aos cursos de graduação e dos projetos extensionistas de caráter tecnológico e competitivo, como Baja, Gravidade Zero e AeroDesign, que se configuram como espaços privilegiados para a inserção dos estudantes em ações interdisciplinares de inovação, desenvolvimento tecnológico e extensão universitária. Esses ambientes favorecem a aprendizagem baseada em projetos, o trabalho em equipe, a resolução de problemas complexos e o desenvolvimento de competências técnicas, empreendedoras e socioemocionais.

A organização curricular do curso favorece a integração entre unidades curriculares e ações de inovação e empreendedorismo por meio de disciplinas como Projetos Extensionistas e Oficinas de Integração, que permitem a inserção dos discentes em contextos reais da sociedade, do setor produtivo e do poder público, para o desenvolvimento de soluções tecnológicas orientadas a demandas concretas. Essas atividades contribuem para a articulação entre teoria e prática, para a consolidação da formação profissional e para o fortalecimento da relação universidade–sociedade.

De forma transversal, o curso incorpora princípios de sustentabilidade \cite{UTFPR_SUST_2022}, economia circular, inovação responsável e fundamentos ambientais, sociais e de governança (ESG) aplicados à tomada de decisão em contextos tecnológicos e organizacionais. Essas práticas são conduzidas por docentes com experiência acadêmica e profissional diversificada, resultando na participação dos estudantes em projetos de pesquisa aplicada, iniciação tecnológica, atividades de extensão, estágios supervisionados, Trabalhos de Conclusão de Curso e na produção e divulgação de resultados técnicos, científicos e tecnológicos.

Dessa forma, a articulação com as políticas institucionais de inovação, empreendedorismo e desenvolvimento tecnológico, ambiental e social consolida-se como elemento estruturante da formação proposta pelo curso de CDIA. Essa articulação contribui para a formação de profissionais capazes de conceber, avaliar e implementar soluções tecnológicas inovadoras, éticas e sustentáveis, comprometidas com o desenvolvimento humano, social e ambiental do Paraná e do Brasil.


\section{Estrutura Curricular e seus componentes}

A estrutura curricular do curso de CDIA do campus Londrina foi concebida de modo a assegurar plena coerência entre o perfil do egresso, as competências profissionais previstas, os objetivos formativos do curso e as demandas contemporâneas do mundo do trabalho, em especial aquelas associadas às transformações digitais, à ciência de dados e à inteligência artificial.

A matriz curricular está organizada de forma progressiva, integrada e interunidade curricular, articulando conhecimentos básicos, específicos e aplicados, com predominância de atividades práticas e contextualizadas, características da formação tecnológica da UTFPR. Os componentes curriculares distribuem-se em períodos letivos que favorecem a evolução gradual da complexidade dos conteúdos e das competências, possibilitando ao estudante construir, aplicar e integrar saberes ao longo de todo o percurso formativo.

A organização curricular contempla unidades curriculares teóricas e práticas, atividades laboratoriais, projetos integradores, atividades extensionistas curricularizadas, estágio curricular supervisionado e Trabalho de Conclusão de Curso (TCC). Essa composição assegura a integração entre teoria e prática, a indissociabilidade entre ensino, pesquisa e extensão, bem como a vivência dos estudantes em situações reais e simuladas do ambiente profissional, favorecendo o desenvolvimento de competências técnicas, científicas e socioemocionais.

A carga horária total do curso, assim como sua distribuição entre componentes curriculares obrigatórios e optativos, atividades práticas, extensão e estágio curricular supervisionado, atende à legislação educacional vigente e às normativas institucionais da UTFPR. A matriz curricular evidencia equilíbrio entre formação técnica específica em Ciência de Dados e Inteligência Artificial, fundamentos científicos e matemáticos, e o desenvolvimento de competências transversais, tais como pensamento crítico, ética profissional, comunicação, trabalho em equipe e uso responsável das tecnologias digitais.

A estrutura curricular prevê flexibilidade formativa por meio da oferta de componentes curriculares optativos, atividades complementares e ações extensionistas, permitindo a adaptação do percurso acadêmico às vocações regionais do campus, às áreas de interesse dos estudantes e às transformações contínuas do campo profissional. Essa flexibilidade também favorece a atualização permanente do currículo, considerando demandas emergentes do setor produtivo, da pesquisa científica e da sociedade.

Os componentes curriculares encontram-se explicitamente articulados às competências e habilidades definidas no perfil do egresso, com objetivos de aprendizagem claramente estabelecidos e coerentes com as metodologias de ensino-aprendizagem adotadas e com os processos de avaliação. A estrutura curricular é periodicamente analisada e atualizada pelo Núcleo Docente Estruturante (NDE) e pelo Colegiado de Curso, assegurando sua aderência às diretrizes institucionais, às normativas nacionais e às necessidades acadêmicas e profissionais emergentes.

Dessa forma, a estrutura curricular do curso de CDIA configura-se como um arranjo formativo consistente, atualizado e socialmente referenciado, que assegura a formação de profissionais tecnológicos qualificados, capazes de atuar de maneira crítica, ética e inovadora em contextos profissionais dinâmicos, complexos e intensivos em tecnologia.

\subsection{Qualificação profissional e empregabilidade}

A organização do curso de Bacharelado CDIA da Universidade Tecnológica Federal do Paraná (UTFPR) – campus Londrina está alinhada à Classificação Internacional Normalizada da Educação Adaptada para Cursos de Graduação (CINE Brasil), às Diretrizes Curriculares Nacionais e às normativas institucionais vigentes, assegurando coerência entre a estrutura curricular, o perfil do egresso, os objetivos formativos e as demandas contemporâneas do mundo do trabalho.

A estrutura do curso é constituída por unidades curriculares organizadas de forma progressiva e integrada, distribuídas ao longo dos períodos letivos, favorecendo o desenvolvimento gradual das competências profissionais requeridas para a atuação do bacharel em CDIA. Os conteúdos, práticas pedagógicas e projetos formativos dialogam diretamente com campos de atuação profissional consolidados e emergentes, assegurando aderência às exigências do setor produtivo, às transformações tecnológicas e à legislação pertinente.

A organização curricular prevê flexibilidade formativa nas dimensões vertical e horizontal, possibilitando melhor aproveitamento acadêmico e a construção de diferentes trajetórias formativas pelos estudantes. A flexibilização vertical compreende a organização das unidades curriculares ao longo dos semestres, contemplando o núcleo de formação específica e a formação não específica, materializada, entre outras possibilidades, pelas unidades curriculares optativas de Humanidades e pelas unidades curriculares optativas e eletivas. Essa organização favorece o aprofundamento progressivo de conhecimentos e competências ao longo do curso.

A flexibilização horizontal baseia-se na ampliação do conceito de currículo, reconhecendo que diversas atividades acadêmicas podem ser consideradas para fins de integralização curricular. Nesse contexto, são consideradas as atividades extraunidade curricular, nas quais se inserem as atividades complementares e as atividades extensionistas curricularizadas. Essas experiências ampliam o contato do estudante com o mundo do trabalho, com a sociedade e com ambientes reais de atuação profissional, fortalecendo competências técnicas, empreendedoras e socioemocionais.

Embora o curso de CDIA não preveja certificações intermediárias formais, uma vez que prioriza uma formação tecnológica integrada e a consolidação das competências profissionais na conclusão do curso, sua estrutura curricular está organizada de modo a promover a qualificação progressiva do estudante. Essa qualificação ocorre por meio de práticas profissionais supervisionadas, projetos integradores, atividades extensionistas, estágios curriculares supervisionados e do Trabalho de Conclusão de Curso, assegurando o desenvolvimento contínuo de competências reconhecidas e valorizadas pelo mercado de trabalho ao longo de todo o percurso formativo.

A qualificação profissional promovida pelo curso é fortalecida pela articulação com organizações públicas e privadas, arranjos produtivos locais e regionais, ambientes de inovação, incubadoras, parques tecnológicos e demais atores do ecossistema de Ciência, Tecnologia e Inovação. Essa integração favorece a inserção qualificada dos estudantes e egressos em atividades profissionais, bem como o desenvolvimento de soluções tecnológicas alinhadas às demandas reais da sociedade e do setor produtivo.

O acompanhamento sistemático da inserção profissional e da trajetória dos egressos subsidia os processos de avaliação e atualização do Projeto Pedagógico do Curso, contribuindo para o fortalecimento da empregabilidade, para a aderência contínua às demandas do mundo do trabalho e para a coerência entre formação acadêmica e exercício profissional.

Dessa forma, a organização do curso assegura a formação de profissionais tecnologicamente qualificados, aptos a atuar em posições compatíveis com o perfil do egresso, dotados de competências técnicas, éticas, críticas e socioemocionais, alinhadas às transformações contemporâneas do mundo do trabalho e às necessidades do desenvolvimento regional e nacional.

\subsection{Matriz curricular}

A Figura \ref{fig:matriz} apresenta a matriz curricular do curso de Bacharelado em Ciência de Dados e Inteligência Artificial, organizada em ciclos formativos articulados, os quais materializam o perfil do egresso e as competências profissionais definidas neste Projeto Pedagógico de Curso. Estão listadas as unidades curriculares usando os acrônimos “BAS” para unidades básicas, “HUM” para unidades básicas humanísticas e “PRO” para as unidades profissionais, evidenciando a progressão formativa e a integração entre fundamentos científicos, conteúdos tecnológicos e aplicações profissionais.

As unidades curriculares também são diferenciadas por coloração, indicando a área de conhecimento predominante em cada unidade curricular. Além disso, cada unidade curricular detalha sua carga horária total, presencial, ensino a distância (EaD) e extensionista. Para auxiliar na compreensão da matriz curricular, a \ref{fig:matriz-legenda} mostra como o significado de cada campo das unidades curriculares.

As aulas são ofertadas no período noturno, sendo a matrícula realizada por unidade curricular, observando os pré-requisitos e a periodização definida no PPC. Esses blocos configuram módulos formativos, os quais favorecem a integração de conteúdos, a articulação teoria–prática e o desenvolvimento de competências em contextos formativos mais complexos.

As temáticas transversais e os fundamentos metodológicos que orientam a organização curricular encontram-se detalhados na Seção \ref{sec:metodologia} – Concepção Metodológica, assegurando coerência entre a matriz curricular, as estratégias pedagógicas e os processos avaliativos adotados pelo curso.

\begin{landscape}
    \begin{figure}[ht]
        \centering
        \caption{Matriz curricular do curso de Bacharelado em Ciência de Dados e Inteligência Artificial.}
        \begin{minipage}{\linewidth}
            \centering 
            
            \includegraphics[width=0.90\linewidth]{figuras/Matriz_Grafica.pdf}

            \ABNTEXfontereduzida \textbf{Fonte:} Autoria própria (2026).
        \end{minipage}
        \label{fig:matriz}
    \end{figure}
\end{landscape}

\begin{figure}[ht]
        \centering
        \caption{Leganda utilizada para descrever as unidades curriculares na matriz.}
        \begin{minipage}{\linewidth}
            \centering 
            
            \includegraphics[width=0.80\linewidth]{figuras/legenda-matriz.png}

            \ABNTEXfontereduzida \textbf{Fonte:} Autoria própria (2026).
        \end{minipage}
        \label{fig:matriz-legenda}
    \end{figure}

\subsubsection{Conteúdos Curriculares}

Cada período letivo da matriz curricular apresentada na Figura \ref{fig:matriz} é detalhado em quadros específicas (Quadros \ref{quadro2.1} a \ref{quadro2.8}), nas quais estão descritas as unidades curriculares correspondentes, identificadas por ciclo formativo e área de conhecimento. Estes quadros explicitam a carga horária presencial, a carga horária mediada por tecnologias digitais (distinguindo atividades síncronas, síncronas mediadas e assíncronas), a carga horária extensionista, a carga horária destinada a conteúdos de formação humanística e a carga horária total de cada período.

Vale ressaltar que a carga horária total nos Quadros \ref{quadro2.1} a \ref{quadro2.8} são calculadas somando as cargas horárias presencial, EaD síncrona, EaD síncrona mediada (com controle de frequência) e EaD assíncrona. As colunas com a carga horária de extensão e de humanidades estão inclusas na carga horária total da unidade curricular.

A descrição das ementas das unidades curriculares encontra-se apresentada nos Apêndices \ref{apendices:obrigatorias} (Quadros das Unidades curriculares obrigatórias) e \ref{apendices:optativas} (Quadros das Unidades curriculares optativas).

As unidades curriculares optativas de formação humanística, podem ser cursadas a partir do segundo período, conforme o interesse acadêmico do estudante e as orientações do curso. As unidades curriculares optativas específicas podem ser cursadas a partir do sexto período. Na matriz curricular todas as disciplinas optativas (humanísticas e específicas) constam no oitavo período (Quadro \ref{quadro2.8}), mas os discentes poderão cursá-las no momento que julgarem adequado, desde que haja compatibilidade entre o período do discente, o horário de oferta da disciplina e a grade horária do discente.

O Quadro \ref{quadro:2.9} apresenta a síntese das cargas horárias de todos os componentes curriculares, evidenciando o atendimento integral às exigências legais e normativas, bem como a carga horária total do curso, em conformidade com as DCN’s e as normativas institucionais da UTFPR.
Conforme as Diretrizes Curriculares da UTFPR os cursos devem ter um Ciclo de Humanidades com carga horária de no mínimo 10\% do total da carga horária destinada às unidades curriculares. A matriz curricular proposta possui 2.595 horas destinadas às unidades curriculares, valor este obtido subtraindo da carga horária total do curso (3.210 horas) as cargas horárias de Estágio Curricular Obrigatório (405 horas), TCC 1 (30 horas), TCC 2 (60 horas) e Atividades Complementares (120 horas). Como o curso oferta um total de 270 horas de disciplinas no Ciclo de Humanidades, considerando optativas e obrigatórias, a matriz contém 10,4\% de carga horária em unidades curriculares de humanidades e atende às regulamentações da UTFPR.
A matriz proposta também atende o mínimo de 10\% da carga horária destinada à extensão, uma vez que há 360 horas de disciplinas extensionistas no curso, o que corresponde a 11,21\% da carga horária total do curso.
Por fim, a matriz curricular possui um total de 555 horas na modalidade de EaD, o que corresponde a 17,29\% da carga horária total do curso. Dessa maneira, a matriz curricular proposta respeita o limite máximo de 30\% de carga horária em EaD permitida para cursos ofertados na modalidade presencial.


% --- QUADRO 2.1 ---
\begin{quadro}[ht!]
\centering
\caption{Lista de Unidades Curriculares do 1º Período}
\label{quadro2.1}
\small

\begin{tabular}{|p{4.2cm}|p{1.2cm}|p{1.2cm}|p{1.2cm}|p{1.2cm}|p{1.2cm}|p{1.2cm}|p{1.2cm}|}
\hline
\rowcolor{gray!10} \textbf{Disciplinas do 1º Período} & \textbf{CH Total} & \textbf{CH Presencial} & \textbf{CH EaD Síncrona} & \textbf{CH EaD Síncrona Mediada} & \textbf{CH EaD Assíncrona} & \textbf{CH Extensão} & \textbf{CH Humanidades} \\ \hline
Algoritmos 1 & 90 & 60 & 0 & 0 & 30 & 0 & 0 \\ \hline
Interpretação e Produção de Textos Científicos e Profissionais & 30 & 30 & 0 & 0 & 0 & 0 & 30 \\ \hline
Introdução à Ciência de Dados e Inteligência Artificial & 60 & 60 & 0 & 0 & 0 & 0 & 0 \\ \hline
Matemática Discreta & 60 & 60 & 0 & 0 & 0 & 0 & 0 \\ \hline
Projeto Extensionista & 120 & 120 & 0 & 0 & 0 & 120 & 0 \\ \hline
\rowcolor{gray!20} \textbf{TOTAL} & \textbf{360} & \textbf{330} & \textbf{0} & \textbf{0} & \textbf{30} & \textbf{120} & \textbf{30} \\ \hline
\end{tabular}
\par\vspace{2pt}
\footnotesize{Fonte: Autoria própria (2026)}
\end{quadro}

% --- QUADRO 2.2 ---
\begin{quadro}[ht!]
\centering
\caption{Lista de Unidades Curriculares do 2º Período}
\label{quadro2.2}
\small
\begin{tabular}{|p{4.2cm}|p{1.2cm}|p{1.2cm}|p{1.2cm}|p{1.2cm}|p{1.2cm}|p{1.2cm}|p{1.2cm}|}
\hline
\rowcolor{gray!10} \textbf{Disciplinas do 2º Período} & \textbf{CH Total} & \textbf{CH Presencial} & \textbf{CH EaD Síncrona} & \textbf{CH EaD Síncrona Mediada} & \textbf{CH EaD Assíncrona} & \textbf{CH Extensão} & \textbf{CH Humanidades} \\ \hline
Álgebra Linear & 45 & 45 & 0 & 0 & 0 & 0 & 0 \\ \hline
Algoritmos 2 & 90 & 60 & 0 & 0 & 30 & 0 & 0 \\ \hline
Fundamentos de Inteligência Artificial & 60 & 45 & 0 & 0 & 15 & 0 & 0 \\ \hline
Gestão Empresarial & 30 & 0 & 0 & 30 & 0 & 0 & 30 \\ \hline
Probabilidade e Estatística & 60 & 60 & 0 & 0 & 0 & 0 & 0 \\ \hline
Programação Orientada a Objetos & 60 & 30 & 0 & 0 & 30 & 0 & 0 \\ \hline
\rowcolor{gray!20} \textbf{TOTAL} & \textbf{345} & \textbf{240} & \textbf{0} & \textbf{30} & \textbf{75} & \textbf{0} & \textbf{30} \\ \hline
\end{tabular}
\par\vspace{2pt}
\footnotesize{Fonte: Autoria própria (2026)}
\end{quadro}

% --- QUADRO 2.3 ---
\begin{quadro}[ht!]
\centering
\caption{Lista de Unidades Curriculares do 3º Período}
\label{quadro2.3}
\small
\begin{tabular}{|p{4.2cm}|p{1.2cm}|p{1.2cm}|p{1.2cm}|p{1.2cm}|p{1.2cm}|p{1.2cm}|p{1.2cm}|}
\hline
\rowcolor{gray!10} \textbf{Disciplinas do 3º Período} & \textbf{CH Total} & \textbf{CH Presencial} & \textbf{CH EaD Síncrona} & \textbf{CH EaD Síncrona Mediada} & \textbf{CH EaD Assíncrona} & \textbf{CH Extensão} & \textbf{CH Humanidades} \\ \hline
Análise Estatística & 60 & 60 & 0 & 0 & 0 & 0 & 0 \\ \hline
Bancos de Dados Relacionais & 60 & 60 & 0 & 0 & 0 & 0 & 0 \\ \hline
Engenharia de Dados & 60 & 60 & 0 & 0 & 0 & 0 & 0 \\ \hline
Estruturas de Dados 1 & 60 & 60 & 0 & 0 & 0 & 0 & 0 \\ \hline
Gerenciamento de Projetos e Desenvolvimento Ágil & 60 & 0 & 0 & 30 & 30 & 0 & 60 \\ \hline
\rowcolor{gray!20} \textbf{TOTAL} & \textbf{300} & \textbf{240} & \textbf{0} & \textbf{30} & \textbf{30} & \textbf{0} & \textbf{60} \\ \hline
\end{tabular}
\par\vspace{2pt}
\footnotesize{Fonte: Autoria própria (2026)}
\end{quadro}

% --- QUADRO 2.4 ---
\begin{quadro}[ht!]
\centering
\caption{Lista de Unidades Curriculares do 4º Período}
\label{quadro2.4}
\small
\begin{tabular}{|p{4.2cm}|p{1.2cm}|p{1.2cm}|p{1.2cm}|p{1.2cm}|p{1.2cm}|p{1.2cm}|p{1.2cm}|}
\hline
\rowcolor{gray!10} \textbf{Disciplinas do 4º Período} & \textbf{CH Total} & \textbf{CH Presencial} & \textbf{CH EaD Síncrona} & \textbf{CH EaD Síncrona Mediada} & \textbf{CH EaD Assíncrona} & \textbf{CH Extensão} & \textbf{CH Humanidades} \\ \hline
Aprendizado de Máquina 1 & 90 & 60 & 0 & 0 & 30 & 0 & 0 \\ \hline
Bancos de Dados Não Relacionais & 60 & 30 & 0 & 0 & 30 & 0 & 0 \\ \hline
Cálculo e Otimização & 60 & 60 & 0 & 0 & 0 & 0 & 0 \\ \hline
Estrutura de Dados 2 & 60 & 60 & 0 & 0 & 0 & 0 & 0 \\ \hline
Sistemas de Computação e Comunicação de Dados & 60 & 30 & 0 & 0 & 30 & 0 & 0 \\ \hline
\rowcolor{gray!20} \textbf{TOTAL} & \textbf{330} & \textbf{240} & \textbf{0} & \textbf{0} & \textbf{90} & \textbf{0} & \textbf{0} \\ \hline
\end{tabular}
\par\vspace{2pt}
\footnotesize{Fonte: Autoria própria (2026)}
\end{quadro}

% --- QUADRO 2.5 ---
\begin{quadro}[ht!]
\centering
\caption{Lista de Unidades Curriculares do 5º Período}
\label{quadro2.5}
\small
\begin{tabular}{|p{4.2cm}|p{1.2cm}|p{1.2cm}|p{1.2cm}|p{1.2cm}|p{1.2cm}|p{1.2cm}|p{1.2cm}|}
\hline
\rowcolor{gray!10} \textbf{Disciplinas do 5º Período} & \textbf{CH Total} & \textbf{CH Presencial} & \textbf{CH EaD Síncrona} & \textbf{CH EaD Síncrona Mediada} & \textbf{CH EaD Assíncrona} & \textbf{CH Extensão} & \textbf{CH Humanidades} \\ \hline
Aprendizado de Máquina 2 & 90 & 60 & 0 & 0 & 30 & 0 & 0 \\ \hline
Comunicação Oral e Estratégica & 30 & 30 & 0 & 0 & 0 & 0 & 30 \\ \hline
Engenharia de Software para Ciência de Dados e IA & 60 & 30 & 0 & 0 & 30 & 0 & 0 \\ \hline
Metodologia de Pesquisa & 30 & 30 & 0 & 0 & 0 & 0 & 30 \\ \hline
Visualização de Dados & 60 & 60 & 0 & 0 & 0 & 0 & 0 \\ \hline
Web API e Integração de Dados & 60 & 30 & 0 & 0 & 30 & 0 & 0 \\ \hline
\rowcolor{gray!20} \textbf{TOTAL} & \textbf{330} & \textbf{240} & \textbf{0} & \textbf{0} & \textbf{90} & \textbf{0} & \textbf{60} \\ \hline
\end{tabular}
\par\vspace{2pt}
\footnotesize{Fonte: Autoria própria (2026)}
\end{quadro}

% --- QUADRO 2.6 ---
\begin{quadro}[ht!]
\centering
\caption{Lista de Unidades Curriculares do 6º Período}
\label{quadro2.6}
\small
\begin{tabular}{|p{4.2cm}|p{1.2cm}|p{1.2cm}|p{1.2cm}|p{1.2cm}|p{1.2cm}|p{1.2cm}|p{1.2cm}|}
\hline
\rowcolor{gray!10} \textbf{Disciplinas do 6º Período} & \textbf{CH Total} & \textbf{CH Presencial} & \textbf{CH EaD Síncrona} & \textbf{CH EaD Síncrona Mediada} & \textbf{CH EaD Assíncrona} & \textbf{CH Extensão} & \textbf{CH Humanidades} \\ \hline
Aprendizado Profundo & 60 & 60 & 0 & 0 & 0 & 0 & 0 \\ \hline
Big Data e Computação em Nuvem & 60 & 30 & 0 & 0 & 30 & 0 & 0 \\ \hline
Desenvolvimento de Novos Negócios & 30 & 0 & 0 & 30 & 0 & 0 & 30 \\ \hline
Implantação e Operação de Modelos & 60 & 60 & 0 & 0 & 0 & 0 & 0 \\ \hline
Oficina de Integração 1 & 120 & 120 & 0 & 0 & 0 & 120 & 0 \\ \hline
Processamento de Linguagem Natural & 60 & 60 & 0 & 0 & 0 & 0 & 0 \\ \hline
\rowcolor{gray!20} \textbf{TOTAL} & \textbf{390} & \textbf{330} & \textbf{0} & \textbf{30} & \textbf{30} & \textbf{120} & \textbf{30} \\ \hline
\end{tabular}
\par\vspace{2pt}
\footnotesize{Fonte: Autoria própria (2026)}
\end{quadro}

% --- QUADRO 2.7 ---
\begin{quadro}[ht!]
\centering
\caption{Lista de Unidades Curriculares do 7º Período}
\label{quadro2.7}
\small
\begin{tabular}{|p{4.2cm}|p{1.2cm}|p{1.2cm}|p{1.2cm}|p{1.2cm}|p{1.2cm}|p{1.2cm}|p{1.2cm}|}
\hline
\rowcolor{gray!10} \textbf{Disciplinas do 7º Período} & \textbf{CH Total} & \textbf{CH Presencial} & \textbf{CH EaD Síncrona} & \textbf{CH EaD Síncrona Mediada} & \textbf{CH EaD Assíncrona} & \textbf{CH Extensão} & \textbf{CH Humanidades} \\ \hline
Computação de Alto Desempenho & 60 & 60 & 0 & 0 & 0 & 0 & 0 \\ \hline
Inteligência Artificial Generativa & 60 & 0 & 0 & 30 & 30 & 0 & 0 \\ \hline
Oficina de Integração 2 & 120 & 120 & 0 & 0 & 0 & 120 & 0 \\ \hline
Sistemas Multiagentes & 60 & 60 & 0 & 0 & 0 & 0 & 0 \\ \hline
Trabalho de Conclusão de Curso 1 & 30 & 0 & 0 & 30 & 0 & 0 & 0 \\ \hline
Visão Computacional & 60 & 60 & 0 & 0 & 0 & 0 & 0 \\ \hline
\rowcolor{gray!20} \textbf{TOTAL} & \textbf{390} & \textbf{300} & \textbf{0} & \textbf{60} & \textbf{30} & \textbf{120} & \textbf{0} \\ \hline
\end{tabular}
\par\vspace{2pt}
\footnotesize{Fonte: Autoria própria (2026)}
\end{quadro}

% --- QUADRO 2.8 ---
\begin{quadro}[ht!]
\centering
\caption{Lista de Unidades Curriculares do 8º Período}
\label{quadro2.8}
\small
\begin{tabular}{|p{4.2cm}|p{1.2cm}|p{1.2cm}|p{1.2cm}|p{1.2cm}|p{1.2cm}|p{1.2cm}|p{1.2cm}|}
\hline
\rowcolor{gray!10} \textbf{Disciplinas do 8º Período} & \textbf{CH Total} & \textbf{CH Presencial} & \textbf{CH EaD Síncrona} & \textbf{CH EaD Síncrona Mediada} & \textbf{CH EaD Assíncrona} & \textbf{CH Extensão} & \textbf{CH Humanidades} \\ \hline
Ética e Legislação para Ciência de Dados e IA & 30 & 0 & 0 & 30 & 0 & 0 & 30 \\ \hline
\rowcolor{gray!20} \textbf{TOTAL} & \textbf{30} & \textbf{0} & \textbf{0} & \textbf{30} & \textbf{0} & \textbf{0} & \textbf{30} \\ \hline
\end{tabular}
\par\vspace{2pt}
\footnotesize{Fonte: Autoria própria (2026)}
\end{quadro}

% --- QUADRO 2.9 ---
\begin{quadro}[htb]
    \centering
    \caption{Consolidação da carga horária total do curso}
    \label{quadro:2.9}
    \begin{tabular}{|l|c|}
        \hline
        \rowcolor{gray!10} \textbf{CONSOLIDAÇÃO DAS CARGAS HORÁRIAS} & \textbf{VALOR (h)} \\ \hline
        CARGA HORÁRIA DISCIPLINAS OBRIGATÓRIAS & 2475 \\ \hline
        CARGA HORÁRIA EAD SÍNCRONA & 0 \\ \hline
        CARGA HORÁRIA EAD SÍNCRONA MEDIADA & 180 \\ \hline
        CARGA HORÁRIA EAD ASSÍNCRONA & 375 \\ \hline
        CARGA HORÁRIA TOTAL EAD & 555 \\ \hline
        CARGA HORÁRIA EM TCC 2 - COMPONENTE CURRICULAR & 60 \\ \hline
        CARGA HORÁRIA DISCIPLINAS OPTATIVAS (NÃO HUMANIDADES) & 120 \\ \hline
        CARGA HORÁRIA DISCIPLINAS ELETIVAS & 0 \\ \hline
        CARGA HORÁRIA ATIVIDADES COMPLEMENTARES & 120 \\ \hline
        CARGA HORÁRIA ESTÁGIO & 405 \\ \hline
        CARGA HORÁRIA DE HUMANIDADES EM DISCIPLINAS OBRIGATÓRIAS & 240 \\ \hline
        CARGA HORÁRIA DE HUMANIDADES EM DISCIPLINAS OPTATIVAS & 30 \\ \hline
        CARGA HORÁRIA TOTAL DE HUMANIDADES & 270 \\ \hline
        CARGA HORÁRIA DE EXTENSÃO EM DISCIPLINAS OBRIGATÓRIAS & 360 \\ \hline
        CARGA HORÁRIA DE EXTENSÃO EM COMPONENTE CURRICULAR & 0 \\ \hline
        CARGA HORÁRIA TOTAL DE EXTENSÃO & 360 \\ \hline
        \rowcolor{gray!20} \textbf{CARGA HORÁRIA TOTAL DO CURSO} & \textbf{3.210} \\ \hline
    \end{tabular}
    \par\vspace{2pt}
    \footnotesize{Fonte: Autoria própria (2026)}
\end{quadro}

\subsection{Estágio Curricular Supervisionado}

O Estágio Curricular Supervisionado constitui ato educativo escolar desenvolvido em ambiente de trabalho, sem geração de vínculo empregatício, e integra o itinerário formativo do estudante como componente curricular previsto no Projeto Pedagógico do Curso (PPC). Sua finalidade é complementar o processo de ensino-aprendizagem, promovendo a articulação entre teoria e prática e preparando o estudante para o exercício profissional ético, crítico e socialmente responsável.

No âmbito legal, destaca-se a Lei nº 11.788, de 25 de setembro de 2008 (Lei do Estágio), que define o estágio como:
\begin{quote}
    “(...) ato educativo escolar supervisionado, desenvolvido no ambiente de trabalho, que visa à preparação para o trabalho produtivo do estudante, proporcionando aprendizagem social, profissional e cultural, por meio de sua participação em atividades de trabalho vinculadas à sua área de formação acadêmico-profissional” \cite{LDE_2008}.
\end{quote}

No plano institucional, o Plano de Desenvolvimento Institucional (PDI) da UTFPR (2018--2022; 2023--2027) \cite{UTFPR_PDI_2017, UTFPR_PDI_2023} estabelece o estágio curricular supervisionado como componente obrigatório para os cursos técnicos e de graduação, reafirmando seu papel na formação profissional. Entre seus objetivos, destacam-se:

\begin{itemize}
    \item Favorecer a inserção qualificada do estudante no mundo do trabalho;
    \item Promover sua adaptação social e profissional às futuras atividades laborais.
\end{itemize}

As Diretrizes para os Cursos de Graduação Regulares da UTFPR (Resolução COGEP nº 142/2022) \cite{UTFPR_DC_2022}, em consonância com as Diretrizes Curriculares Nacionais \cite{DCN_2011, MEC_2016} e com as recomendações da Sociedade Brasileira de Computação (SBC) \cite{SBC_2017, SBC_REF_CD_2023, SBC_REF_IA_2024}, determinam que os cursos de bacharelado devem prever Estágio Curricular Obrigatório com carga horária mínima de 360 horas.

O Regulamento dos Estágios Curriculares Supervisionados da UTFPR (Resolução Conjunta nº 01, de 02 de junho de 2020) \cite{Estagio_2020} estabelece que:

\begin{itemize}
    \item O estágio integra o PPC e o percurso formativo do estudante;
    \item Visa ao desenvolvimento de competências profissionais e à contextualização curricular;
    \item Deve ser realizado em áreas compatíveis com o perfil profissional definido no curso.
\end{itemize}

A UTFPR, por meio do Departamento de Estágios e Empregos (DEPEC) e da Divisão de Estágios (DIEEM), opera um Sistema de Estágios integrado em cada campus, permitindo:

\begin{itemize}
    \item O cadastro e a validação de empresas e instituições para oferta de estágios e/ou oportunidades de trabalho;
    \item O acompanhamento sistemático e a supervisão das atividades desenvolvidas pelos estudantes.
\end{itemize}

Esse sistema assegura controle acadêmico e administrativo, conformidade legal e registro das experiências formativas, constituindo mecanismo estruturado de retroalimentação da formação profissional.

Atualmente, a instituição mantém convênio com mais de 900 unidades concedentes, número continuamente ampliado mediante validação das coordenações de curso. O processo de qualificação das vagas assegura a coerência entre as atividades propostas e a área de formação do estudante, garantindo aderência ao perfil profissional definido no PPC. Os convênios firmados contemplam dispositivos que:
\begin{itemize}
    \item Contribuem para a atualização contínua do PPC, por meio de reuniões com supervisores das empresas e análises de aderência entre atividades profissionais e o perfil do egresso;
    \item Fortalecem o desenvolvimento das competências previstas, mediante oferta de vagas qualificadas, projetos conjuntos de inovação, participação de profissionais externos em bancas, palestras técnicas e avaliações situacionais;
    \item Subsidiam estudos de empregabilidade e acompanhamento de egressos, com base em dados coletados junto às organizações parceiras;
    \item Evidenciam a relação estruturada entre curso e setor produtivo, expressa em contrapartidas como oferta de vagas, apoio a atividades formativas e colaboração em ações de extensão tecnológica.
\end{itemize}

No curso de curso de CDIA, o estudante poderá realizar duas modalidades de estágio: Estágio Não Obrigatório e Estágio Curricular Obrigatório.
O Estágio Não Obrigatório é de caráter opcional e formativo complementar, poderá ser iniciado a partir do 2º período letivo. Tem como finalidade proporcionar inserção gradual no ambiente profissional, permitindo ao estudante:
\begin{itemize}
    \item Familiarizar-se com rotinas organizacionais;
    \item Desenvolver competências básicas;
    \item Observar processos produtivos;
    \item Participar de atividades técnicas compatíveis com seu nível de formação.    
\end{itemize}

Seus resultados concentram-se no desenvolvimento de atitudes profissionais, responsabilidade, trabalho em equipe e progressiva autonomia.

O Estágio Curricular Obrigatório constitui componente curricular obrigatório, com carga horária mínima de 405 horas, requisito para integralização do curso e obtenção do diploma, conforme Resolução COGEP nº 142/2022 \cite{UTFPR_DC_2022}. Poderá ser realizado a partir do 6º período do curso.

Caracteriza-se por envolver atividades compatíveis com o nível terminal da formação, demandando aplicação integrada de conhecimentos técnicos, metodológicos e tecnológicos. O estudante deverá demonstrar competências previstas no perfil do egresso, atuando em atividades de maior complexidade, tais como:

\begin{itemize}
    \item Análise e modelagem de dados;
    \item Desenvolvimento e avaliação de modelos de aprendizado de máquina;
    \item Implementação de soluções de engenharia de dados;
    \item Desenvolvimento e avaliação de sistemas inteligentes;
    \item Otimização de processos baseados em dados;
    \item Apoio à tomada de decisão orientada por dados.
\end{itemize}

Em ambas as modalidades, as atividades deverão manter relação direta e comprovada com as áreas de atuação do profissional de CDIA, podendo ocorrer em organizações públicas ou privadas, empresas de tecnologia, indústrias, instituições financeiras, startups, laboratórios de inovação ou iniciativas vinculadas ao ecossistema regional de inovação.

O acompanhamento do estágio é realizado por docente orientador do curso e por supervisor da instituição concedente. O processo é mediado por sistemas institucionais informatizados que permitem:

\begin{itemize}
    \item Registro formal das atividades;
    \item Acompanhamento da carga horária;
    \item Comunicação entre as partes envolvidas;
    \item Avaliação das experiências formativas.
\end{itemize}

O curso contará com Professor Responsável pelas Atividades de Estágio (PRAE), a quem compete avaliar a pertinência das atividades, analisar a documentação (plano de estágio, relatórios e avaliações), acompanhar o desenvolvimento das atividades e proceder ao lançamento da nota no sistema acadêmico.

Os relatórios e registros produzidos pelos estudantes são analisados periodicamente pelo Núcleo Docente Estruturante (NDE) e pelo Colegiado de Curso, contribuindo para o aprimoramento contínuo do PPC e para o fortalecimento da integração com o mundo do trabalho.

A jornada do estágio deverá ser compatível com o horário acadêmico do estudante, não ultrapassando 30 horas semanais durante o período letivo, conforme a Lei nº 11.788/2008 \cite{LDE_2008}, podendo atingir 40 horas semanais no período de férias acadêmicas.

Dessa forma, o Estágio Curricular Supervisionado configura-se como eixo estruturante da formação do Bacharel em CDIA, promovendo a integração entre conhecimentos, habilidades e atitudes; fortalecendo a aprendizagem ao longo da vida; aproximando o estudante de contextos reais e complexos da profissão; e apoiando sua preparação para desafios contemporâneos associados à inovação, à transformação digital e às demandas tecnológicas emergentes.

\subsection{Atividades Complementares}

As Atividades Complementares têm por finalidade enriquecer o processo de ensino-aprendizagem, contribuindo para a complementação da formação social, humana, cultural, científica e profissional do estudante, em consonância com o perfil do egresso e as competências previstas no Projeto Pedagógico do Curso. Caracterizam-se pela flexibilidade quanto à forma de realização e distribuição da carga horária, respeitado o controle do tempo total de dedicação ao longo do curso, conforme orientações do Parecer CNE/CES nº 492/2001.

No âmbito do curso, as Atividades Complementares devem ser integralizadas pelo estudante até o oitavo período, em conformidade com o Regulamento das Atividades Complementares dos Cursos de Graduação da UTFPR \cite{UTFPR_AC_2007}, instituído pela Resolução COGEP nº 179, de 04 de agosto de 2022, que regulamenta as Atividades Complementares (ACs) no âmbito institucional.

As Atividades Complementares estão organizadas em dois grupos, conforme previsto na regulamentação institucional:

\begin{itemize}
    \item \textbf{Grupo 1 – Formação Social, Humana e Cultural:} contempla atividades de cunho cultural, esportivo, social e formativo, que contribuam para o desenvolvimento da cidadania, da sensibilidade cultural e da formação ética do estudante. Incluem-se, entre outras, a participação em eventos culturais e esportivos, projetos solidários e assistencialistas, cursos de línguas estrangeiras, atividades artísticas e culturais.
    \item \textbf{Grupo 2 – Formação Científica e Técnico-Profissional:} contempla atividades voltadas ao desenvolvimento do espírito científico, do pensamento crítico e da qualificação técnica, tais como participação em eventos acadêmicos e científicos (semanas acadêmicas, congressos, seminários, palestras e conferências), cursos de atualização acadêmica e profissional, atividades de iniciação científica, monitoria acadêmica e outras ações correlatas.
\end{itemize}

As quantificações e a forma de comprovação desta carga horária se apresentam em ato normativo próprio. O Professor Responsável pelas Atividades Complementares (PRAC), é a quem compete a análise, validação, registro acadêmico e, quando necessário, a conferência dos documentos originais, em conformidade com os procedimentos institucionais estabelecidos.

O curso incentiva a participação dos estudantes em atividades complementares diversificadas, preferencialmente desenvolvidas fora do ambiente formal de sala de aula, de modo a ampliar repertórios culturais, científicos e profissionais, respeitadas as normas institucionais e os limites de carga horária definidos. O detalhamento das atividades aceitas, os critérios de contabilização de horas e os procedimentos administrativos encontram-se descritos em norma complementar específica do curso, em consonância com a Resolução COGEP nº 179/2022.

Dessa forma, as Atividades Complementares contribuem de maneira efetiva para a formação integral do estudante, estimulando a autonomia, o pensamento reflexivo, a inserção em diferentes contextos socioculturais e o desenvolvimento de competências alinhadas às demandas locais, regionais, nacionais e globais.

\subsection{Trabalho de Conclusão de Curso}

No contexto do curso de CDIA no campus Londrina, o Trabalho de Conclusão de Curso constitui um componente essencial para a consolidação da formação tecnológica, integrando os conhecimentos construídos ao longo do percurso acadêmico e estimulando a autonomia intelectual do estudante. Seu objetivo central é possibilitar que o discente mobilize competências técnicas, científicas e humanísticas para investigar, analisar e propor soluções inovadoras em temas relacionados à Inteligência Artificial, ciência de dados e demais áreas tecnológicas correlatas. Ao desenvolver o TCC, o estudante é convidado a exercitar o pensamento crítico, a capacidade de resolução de problemas complexos e a responsabilidade ética diante dos impactos sociais, econômicos e culturais das tecnologias digitais.

O desenvolvimento, acompanhamento e avaliação do TCC obedecerão às normas institucionais da UTFPR. As regras específicas de dos possíveis formatos de elaboração, as formas de apresentação, avaliação, orientação e apresentação serão definidas em ato normativo próprio, aprovado pelo colegiado do curso e consolidado nas Normas Complementares de TCC, em consonância com o Regulamento de Trabalho de Conclusão de Curso para os cursos de graduação da UTFPR (Resolução COGEP/UTFPR nº 180/2022) \cite{UTFPR_TCC_2018} ou documento vigente. Esses instrumentos garantem a padronização dos procedimentos, a qualidade acadêmica dos trabalhos e a coerência com o perfil profissional no Bacharelado em CDIA.

Após aprovação, os Trabalhos de Conclusão de Curso são registrados e disponibilizados em repositório institucional ou sistema acadêmico da UTFPR, garantindo sua publicização como recursos educacionais abertos, quando aplicável, respeitados os aspectos éticos, legais e de propriedade intelectual.

A operacionalização do TCC é acompanhada pelo Professor Responsável pelo TCC (PRATCC), com supervisão do NDE e do Colegiado de Curso, que avaliam periodicamente sua aderência ao Projeto Pedagógico, sua contribuição para o desenvolvimento das competências do egresso e sua efetividade como estratégia de integração entre formação acadêmica, prática profissional, inovação e compromisso social. Eventuais situações excepcionais são analisadas pelo Colegiado, mediante justificativa formal, assegurando flexibilidade institucional sem prejuízo da qualidade formativa.

Dessa forma, o Trabalho de Conclusão de Curso consolida-se como componente estruturante da formação no curso de CDIA, contribuindo para a qualificação profissional, a empregabilidade e a formação de tecnólogos capazes de atuar de maneira crítica, inovadora e eticamente responsável no contexto da transformação digital e do desenvolvimento tecnológico contemporâneo.

\subsection{Inserção Curricular na Extensão}

Implementando as políticas institucionais de extensão da Universidade Tecnológica Federal do Paraná (UTFPR), estabelecidas pela Resolução Conjunta COGEP/COEMP/UTFPR nº 2, de 8 de agosto de 2025 \cite{UTFPR_EXT_2022}, o curso operacionaliza a extensão universitária por meio dos eixos formativos de extensão previamente definidos, integrando-os de forma estruturada à sua organização curricular.

Esses eixos orientam o planejamento, a execução e a avaliação das atividades extensionistas inseridas na matriz curricular. Cada eixo está diretamente articulado ao perfil do egresso e às competências profissionais previstas, assegurando coerência entre formação acadêmica, prática profissional e compromisso social.

A implementação da extensão no curso se materializa integralmente por meio de Unidades Curriculares Extensionistas obrigatórias, as quais são sumarizadas no Quadro \ref{quadro:2.10}.

\begin{quadro}[htb]
    \centering
    \caption{Unidades Curriculares de extensão}
    \label{quadro:2.10}
    \begin{tabular}{|c|l|c|}
        \hline
        \textbf{Período} & \textbf{Disciplina} & \textbf{Carga Horária} \\ \hline
        1º & Projeto Extensionista & 120 horas \\ \hline
        6º & Oficina de Integração 1 & 120 horas \\ \hline
        7º & Oficina de Integração 2 & 120 horas \\ \hline
        & TOTAL & 360 horas \\ \hline
    \end{tabular}
    
    \vspace{1ex}
    \small Fonte: Autoria prória (2026)
\end{quadro}

A disciplina de Projeto Extensionista no 1º período do curso tem como objetivo inserir o discente em atividades extensionistas do campus com cunho social e comunitário, visando também divulgar a universidade e seus cursos. As disciplinas de Oficina de Integração 1 e 2 nos períodos mais avançados do curso buscam inserir o discente em projetos de aplicação dos conhecimentos de ciência de dados e inteligência artificial em problemas reais da comunidade externa à universidade. A cada semestre serão levantados problemas relevantes que serão tópico da disciplina, dando aos discentes o papel de protagonista no entendimento dos problemas reais e proposta de soluções, com o auxílio de docentes do curso.

A operacionalização da carga horária extensionista é acompanhada pelo Professor Responsável pela extensão (PRAEXT), com supervisão do NDE e do Colegiado de Curso, que avaliam periodicamente sua contribuição para o desenvolvimento das competências do egresso e sua efetividade como estratégia de integração entre formação acadêmica, prática profissional, inovação e compromisso social. Os resultados das atividades extensionistas implementadas no curso são analisados periodicamente e utilizados como subsídio para o aprimoramento contínuo da estrutura curricular, para o fortalecimento da integração com o mundo do trabalho e para a ampliação dos impactos sociais, econômicos, culturais, ambientais e tecnológicos no território de atuação do curso.

Para o curso de Bacharelado em CDIA do campus Londrina, a extensão assume papel estratégico ao promover a aplicação do conhecimento tecnológico em contextos reais, permitindo que estudantes desenvolvam soluções de impacto social, tecnológico e econômico. Projetos envolvendo inclusão digital, automação de processos, aplicações de IA em saúde, educação, indústria, segurança digital, sustentabilidade e inovação social exemplificam ações que materializam o compromisso da UTFPR com o desenvolvimento regional e com a formação de profissionais capazes de responder aos desafios contemporâneos.

Assim, a curricularização da extensão no curso busca fortalecer a formação integral dos estudantes, ampliar o diálogo com a sociedade, valorizar práticas colaborativas e fomentar experiências que integrem tecnologia, humanismo e responsabilidade social — princípios constituintes da missão institucional da UTFPR. Alinhada ao PDI \cite{UTFPR_PDI_2017, UTFPR_PDI_2023} e às diretrizes nacionais \cite{DCN_2011, MEC_2016}, a extensão universitária contribui para a consolidação de uma formação crítica, interunidade curricular e comprometida com a transformação social e tecnológica, em sintonia com as demandas e potencialidades das áreas STEM e da Inteligência Artificial.

\subsubsection{Integração do curso com a sociedade}

A integração do curso de Bacharelado em CDIA com a sociedade constitui eixo estruturante da proposta formativa e materializa o compromisso social da UTFPR com o desenvolvimento sustentável do território em que está inserida. Essa integração ocorre por meio do desenvolvimento sistemático de ações, programas, projetos, cursos, oficinas, eventos de extensão, hackathons e prestação de serviços tecnológicos, planejados de forma regular e articulados à estrutura curricular do curso.

As ações desenvolvidas pelo curso estão orientadas para a aplicação do conhecimento tecnológico em demandas reais da comunidade, envolvendo organizações públicas, privadas e do terceiro setor, e promovendo impacto social, econômico, cultural, ambiental e tecnológico no contexto local e regional. Essas iniciativas possibilitam ao estudante vivenciar situações autênticas de atuação profissional, fortalecendo o reconhecimento da função social do tecnólogo e contribuindo para a formação ética, crítica e socialmente responsável. No contexto de Londrina, podemos destacar alguns exemplos de como o curso pode se integrar com a sociedade:
\begin{itemize}
    \item Projetos de extensão envolvendo análise de dados públicos, auxiliando órgãos públicos municipais e estaduais na construção de dashboards e relatórios;
    \item Apoio tecnológico a Pequenas e Médias Empresas (PMEs), auxiliando-as no desenvolvimento de soluções que resultem em melhorias do desempenho operacional de processos produtivos;
    \item Participação e organização de eventos locais, como mentorias em hackatons, atuação em projetos de incubação, participação em iniciativas do ecossistema local de inovação;
    \item Inserção por meio de estágios nas empresas de base tecnológica de Londrina;
    \item Realização de cursos de extensão para a comunidade, buscando capacitar pessoas no uso de tecnologias de informação, etc.
\end{itemize}

As atividades de integração com a sociedade são avaliadas periodicamente pelo Núcleo Docente Estruturante (NDE) e acompanhadas pela coordenação e pelo colegiado do curso, considerando sua aderência ao perfil do egresso, às competências profissionais previstas no PPC e às demandas do território. Os resultados dessas ações subsidiam o aprimoramento contínuo do curso, bem como a atualização de práticas pedagógicas e conteúdos curriculares.

Além do impacto formativo, as ações extensionistas e de interação social desenvolvidas pelo curso de CDIA contribuem para a qualificação da pesquisa aplicada, quando pertinente, ao gerar problemas, dados e desafios reais que retroalimentam projetos de investigação, inovação e desenvolvimento tecnológico. Dessa forma, o curso fortalece a indissociabilidade entre ensino, pesquisa e extensão, ampliando sua relevância social e acadêmica.

Assim, a integração do curso com a sociedade consolida-se como prática institucionalizada e permanente, promovendo o desenvolvimento social e sustentável das comunidades atendidas, ao mesmo tempo em que forma profissionais capazes de atuar de maneira qualificada, inovadora e comprometida com as transformações contemporâneas da sociedade.

\subsection{Formação Humanística}

No curso de CDIA a formação tecnológica articula-se à formação integral, reconhecendo que o exercício profissional em Inteligência Artificial exige não apenas domínio técnico, mas também compreensão crítica das dimensões éticas, sociais, culturais e ambientais envolvidas no desenvolvimento científico e tecnológico.

Em conformidade com a Resolução COGEP/UTFPR nº 142/2022 \cite{UTFPR_DC_2022}, o curso de CDIA incorpora em sua matriz curricular um ciclo de humanidades correspondente a, no mínimo 10\% da carga horária destinada às unidades curriculares. O Ciclo de Humanidades é composto por unidades curriculares das áreas de ciências humanas, ciências sociais aplicadas, linguagem e artes, podendo incluir conteúdos de saúde, qualidade de vida e atividades de extensão, assegurando a integração entre a formação técnica e humanista ao longo do percurso formativo. A Quadro \ref{quadro:2.11} apresenta as disciplinas do Ciclo de Humanidades na matriz curricular do curso.


\begin{quadro}[htb]
    \centering
    \caption{Unidades curriculares do Ciclo de Humanidades}
    \label{quadro:2.11}
    \begin{tabular}{|c|l|c|}
        \hline
        \textbf{Período} & \textbf{Nome da Disciplina} & \textbf{CH} \\ \hline
        1 & Interpretação e Produção de Textos Científicos e Profissionais & 30 \\ \hline
        2 & Gestão Empresarial & 30 \\ \hline
        3 & Gerenciamento de Projetos e Desenvolvimento Ágil & 60 \\ \hline
        5 & Comunicação Oral e Estratégica & 30 \\ \hline
        5 & Metodologia de Pesquisa & 30 \\ \hline
        6 & Desenvolvimento de Novos Negócios & 30 \\ \hline
        8 & Ética e Legislação para CDIA & 30 \\ \hline
        8 & Optativa de Humanidades & 30 \\ \hline
        & TOTAL & 270 \\ \hline
    \end{tabular}
    \fonte{Autoria prória (2026)}
\end{quadro}


A concepção da estrutura curricular e de seus componentes propicia a abordagem sistemática em unidade curriculares especializadas e pertinentes às políticas educacionais e sociais, incluindo:

\begin{itemize}
    \item Educação ambiental e sustentabilidade, com foco nos impactos socioambientais das tecnologias digitais e da Inteligência Artificial \cite{Brasil_1997};
    \item Direitos humanos, ética profissional e responsabilidade social no desenvolvimento e uso de sistemas tecnológicos;
    \item Temas relacionados à pessoa com deficiência, acessibilidade digital e tecnologias assistivas \cite{Brasil_2015_Deficiencia, Brasil_2011_Deficiencia};
    \item Relações étnico-raciais e o estudo da história e cultura afro-brasileira, africana e indígena, promovendo uma formação crítica, inclusiva e socialmente referenciada \cite{Brasil_Raciais_2004}.
\end{itemize}

Esses conteúdos são também trabalhados de forma interunidade curricular e transversal, integrando-se às unidades curriculares técnicas, aos projetos integradores e às atividades de extensão, evitando sua fragmentação e assegurando sua articulação com os contextos reais de atuação profissional. Assim, é importante destacar que as disciplinas extensionistas, apesar de não serem computadas no Ciclo de Humanidades, tem como foco o desenvolvimento das temáticas listadas anteriormente. Essa abordagem interunidade curricular também favorece a reflexão sobre inovação, empreendedorismo tecnológico e responsabilidade social, ampliando a compreensão do papel do tecnólogo na sociedade contemporânea.

A estrutura curricular evidencia, ainda, a oferta da unidades curriculars de Libras 1 e Libras 2, em atendimento à legislação vigente, ampliando a formação inclusiva e a sensibilização dos estudantes para a comunicação acessível e para a diversidade humana. As unidades curriculares de Libras constam como optativas da área de Humanidades na matriz curricular para flexibilizar a formação dos estudantes. Por se tratar de disciplina prevista em vários cursos de graduação do campus Londrina, semestralmente são ofertadas mais de uma turma de cada uma dessas disciplinas, garantindo aos estudantes a oportunidade de cursá-las.

No âmbito da integração do curso com a sociedade, os temas transversais de sustentabilidade, diversidade e responsabilidade social estão contemplados e aplicados de forma concreta por meio de ações extensionistas, projetos e atividades desenvolvidas em interação com a comunidade acadêmica e externa, incluindo ambientes profissionais, organizações públicas, privadas e do terceiro setor. Essas ações contribuem para a formação cidadã e ética profissional, ao mesmo tempo em que fortalecem o compromisso social da UTFPR e o impacto positivo do curso no território em que está inserido. Na matriz curricular os conteúdos de educação ambiental e sustentabilidade são explicitamente abordados nas disciplinas extensionistas obrigatórias Projeto Extensionista (1º Período) e Oficina de Integração 2 (7º Período) e na disciplina Ética e Legislação para Ciência de Dados e IA (8º Período). Além disso, os discentes poderão expandir seus conhecimentos nessas áreas por meio das disciplinas optativas “Meio Ambiente e Sociedade” e “Economia do Meio Ambiente”.

Dessa forma, a formação humanística no curso de CDIA consolida-se como elemento estruturante do projeto pedagógico, assegurando que a formação técnica esteja indissociavelmente integrada a uma visão crítica, ética, inclusiva e socialmente responsável da tecnologia, em consonância com as diretrizes institucionais, nacionais e com as demandas contemporâneas da sociedade.

\subsection{Promoção do empreendedorismo e da inovação}

O curso de CDIA do campus Londrina implementa ações e conteúdos voltados à promoção do empreendedorismo e da inovação como componentes estruturantes de sua proposta formativa, em consonância com o perfil do egresso e com as competências profissionais previstas no Projeto Pedagógico do Curso.

As ações relacionadas ao empreendedorismo e à inovação estão inseridas na matriz curricular, por meio das disciplinas de Oficina de Integração 1 e 2, bem como a disciplina de Desenvolvimento de Novos Negócios. Tais ações também se materializam em programas, projetos e ações de extensão coordenados por docentes do curso, possibilitando ao estudante compreender o processo tecnológico, identificar oportunidades e desenvolver soluções inovadoras aplicáveis à área de Ciência de Dados e Inteligência Artificial, nos diversos setores da economia e da sociedade que podem ser beneficiados por tais soluções.

Essas iniciativas fomentam o desenvolvimento da capacidade empreendedora, da criatividade, da autonomia e da tomada de decisão, sendo estimuladas por meio de parcerias com o setor produtivo, incluindo o Arranjo Produtivo Local (APL) de Tecnologia da Informação e Comunicação de Londrina, bem como os demais atores do ecossistema de inovação da cidade. A proximidade com tais entidades contribui para a aproximação do curso com o mundo do trabalho e para a aplicação do conhecimento em contextos reais.

O curso encontra-se articulado a ambientes de inovação, tais como incubadoras, hotéis tecnológicos, parques tecnológicos, núcleos de inovação, ecossistemas locais ou regionais, por meio de convênios, termos de cooperação, projetos institucionais e atividades conjuntas. Essas articulações possibilitam a participação ativa de docentes e discentes em ações de inovação, desenvolvimento tecnológico e empreendedorismo, ampliando as oportunidades de aprendizagem e de inserção profissional.

As ações de empreendedorismo e inovação desenvolvidas no âmbito do curso são acompanhadas, monitoradas e avaliadas pela coordenação do curso e pelo Núcleo Docente Estruturante (NDE), assegurando sua coerência com as diretrizes institucionais, sua continuidade e sua contribuição efetiva para o desenvolvimento das competências profissionais do egresso.

Dessa forma, o curso consolida uma formação tecnológica orientada à inovação, ao empreendedorismo e ao desenvolvimento sustentável, alinhada às demandas contemporâneas da sociedade e do mundo do trabalho.


\subsection{Modalidade EaD para o Curso Presencial}

O curso de CDIA do campus Londrina é ofertado no formato presencial, em conformidade com os princípios institucionais da UTFPR e com as Diretrizes Curriculares da Educação Profissional e Tecnológica \cite{Analise_DNC_2021}. Além disso, o curso observa os seguintes marcos regulatórios:

\begin{itemize}
    \item O Decreto nº 12.456, de 19 de maio de 2025, que estabelece o Novo Marco Regulatório da Educação a Distância no Brasil \cite{Brasil_EaD_2025};
    \item Os Referenciais de Qualidade para a Educação Superior à Distância, do Ministério da Educação/Secretaria de Educação a Distância, de maio de 2025, e;
    \item A Resolução COGEP/UTFPR nº 742, de novembro de 2025, que regulamenta a oferta de cursos de graduação nas modalidades a distância e semipresencial, bem como a oferta de atividades a distância em cursos de graduação presenciais no âmbito da Universidade Tecnológica Federal do Paraná \cite{UTFPR_EAD_2022}.
\end{itemize}

A incorporação de unidades curriculares no formato integral ou parcial a distância no curso de CDIA, dentro do limite legal de 30\% da carga horária total do curso, constitui uma estratégia pedagógica alinhada às transformações contemporâneas da Educação Superior. Nesse contexto, a EaD é planejada, avaliada e integradaa este PPC, orientando-se por princípios de qualidade, ampliação do acesso à educação superior, promoção de espaços de interação e dialogicidade, desenvolvimento de competências e habilidades, diversidade de materiais didáticos e recursos tecnológicos, acessibilidade, inclusão, fortalecimento da identidade institucional e desenvolvimento profissional.

As disciplinas ofertadas no formato a distância são estruturadas e organizadas de modo a articular fundamentos teóricos, metodologias pedagógicas inovadoras, mediação pedagógica qualificada e uso consistente das Tecnologias Digitais de Informação e Comunicação (TDIC). Nesse cenário, os discentes do curso de CDIA assumem papel central em seu processo formativo, desenvolvendo autonomia intelectual, competências digitais, organização do tempo e capacidade de aprendizagem colaborativa, em consonância com os preceitos de uma Educação Superior pautada na flexibilidade curricular, na inovação pedagógica, no uso crítico e ético das TDIC e na indissociabilidade entre ensino, pesquisa e extensão, contribuindo para uma formação integral, crítica e comprometida com as demandas contemporâneas da sociedade e do mundo do trabalho. O Quadro \ref{quadro:2.12} apresenta as unidades curriculares da matriz que serão ofertadas parcialmente ou integralmente à distância no curso.

Observando o Quadro \ref{quadro:2.12} é possível notar que o curso distribui ao longo de todos os períodos a carga horária à distância, permitindo que haja maior flexibilidade em sua definição de horário e em sua permanência presencial na universidade, o que é de grande relevância para cursos no período noturno. Por exemplo, disciplinas como Algoritmos 1 e 2 possuem 60 horas de carga horária presencial para o ensino mais próximo entre discentes e docentes, mas também possuem 30 horas de carga horária à distância para que os estudantes possam complementar seu aprendizado com atividades assíncronas. Há também disciplinas integralmente à distância, como Gerenciamento de Projetos e Desenvolvimento Ágil, na qual a carga horária à distância será subdividida em atividades síncronas mediadas, com a interação em tempo real com o docente, e em atividades assíncronas de autoestudo e realização de atividades e trabalhos.

Com base nas cargas horárias apresentadas no Quadro \ref{quadro:2.12}, o discente terá ao todo 555 horas à distância na matriz curricular, o que corresponde a 17,32\% da carga horária total do curso. Vale observar que, no caso das disciplinas optativas do curso, o discente poderá cursá-las tanto presencialmente como à distância. Caso escolha por cursar todas as disciplinas optativas à distância, sua carga horária em EaD chegará a 705 horas, o que corresponde a 22\% da carga horária total do curso, ainda assim dentro do limite de 30\% para cursos presenciais. Dessa maneira, o discente terá flexibilidade em escolher como cursará as disciplinas optativas, podendo ainda cursá-las em outros campi da UTFPR que ofertam disciplinas à distância, evidenciando a sinergia com os demais cursos $STEM^{IA}$ ofertados na universidade.

\begin{quadro}[htb]
    \centering
    \caption{Unidades curriculares com carga horária EaD}
    \label{quadro:2.12}
    \vspace{0.2cm}
    \resizebox{\textwidth}{!}{ % Ajusta o quadro à largura da página, caso fique muito grande
    \begin{tabular}{|c|l|c|c|c|c|c|c|}
        \hline
        \textbf{Período} & \textbf{Nome da Disciplina} & \textbf{CH Pres} & \textbf{CH Assinc} & \textbf{CH Sinc Med} & \textbf{CH Sinc} & \textbf{CH Dist Total} & \textbf{CH Total} \\ \hline
        1º & Algoritmos 1 & 60 & 30 & 0 & 0 & 30 & 90 \\ \hline
        2º & Algoritmos 2 & 60 & 30 & 0 & 0 & 30 & 90 \\ \hline
        2º & Fundamentos de Inteligência Artificial & 45 & 15 & 0 & 0 & 15 & 60 \\ \hline
        2º & Gestão Empresarial & 0 & 0 & 30 & 0 & 30 & 30 \\ \hline
        2º & Programação Orientada a Objetos & 30 & 30 & 0 & 0 & 30 & 60 \\ \hline
        3º & Gerenciamento de Projetos e Desenvolvimento Ágil & 0 & 30 & 30 & 0 & 60 & 60 \\ \hline
        4º & Aprendizado de Máquina 1 & 60 & 30 & 0 & 0 & 30 & 90 \\ \hline
        4º & Bancos de Dados Não Relacionais & 30 & 30 & 0 & 0 & 30 & 60 \\ \hline
        4º & Sistemas de Computação e Comunicação de Dados & 30 & 30 & 0 & 0 & 30 & 60 \\ \hline
        5º & Aprendizado de Máquina 2 & 60 & 30 & 0 & 0 & 30 & 90 \\ \hline
        5º & Engenharia de Software para IA & 30 & 30 & 0 & 0 & 30 & 60 \\ \hline
        5º & Web API e Integração de Dados & 30 & 30 & 0 & 0 & 30 & 60 \\ \hline
        6º & Big Data e Computação em Nuvem & 30 & 30 & 0 & 0 & 30 & 60 \\ \hline
        6º & Desenvolvimento de Novos Negócios & 0 & 0 & 30 & 0 & 30 & 30 \\ \hline
        7º & Inteligência Artificial Generativa & 0 & 30 & 30 & 0 & 60 & 60 \\ \hline
        7º & TCC 1 & 0 & 0 & 30 & 0 & 30 & 30 \\ \hline
        8º & Ética e Legislação para CDIA & 0 & 0 & 30 & 0 & 30 & 30 \\ \hline
        \multicolumn{2}{|c|}{\textbf{TOTAL}} & \textbf{465} & \textbf{375} & \textbf{180} & \textbf{0} & \textbf{555} & \textbf{1020} \\ \hline
    \end{tabular}
    }
    \fonte{Autoria prória (2026)}
\end{quadro}

Em consonância com os documentos nacionais e institucionais vigentes \cite{UTFPR_EAD_2022, Analise_DNC_2021}, as unidades curriculares ofertadas no formato integral ou parcial a distância observam critérios claros de planejamento didático-pedagógico, carga horária compatível, acompanhamento sistemático por docentes formadores e, quando necessário, por mediadores pedagógicos, além de mecanismos de avaliação contínua. Os ambientes virtuais de aprendizagem (AVAs) e demais recursos tecnológicos são utilizados como espaços de interação, produção do conhecimento e acompanhamento acadêmico, assegurando a efetiva participação dos discentes e a integração entre teoria e prática.

A equipe multiunidade curricular será composta pelos seguintes profissionais: pedagogos, psicólogo, design educacional, design gráfico, revisor, especialista em acessibilidade, coordenador de estágio, coordenador de extensão, dentre outros. Esses profissionais serão responsáveis pelo acompanhamento sistemático das ações pedagógicas do curso, buscando apoiar docentes e discentes durante o processo de ensino-aprendizagem na EaD.

Os materiais didáticos serão diversificados quanto a formatos e linguagens, organizados de forma autoexplicativa e dialógica, de modo a possibilitar que o discente construa seu próprio itinerário formativo. Cada unidade curricular buscará desenvolver uma variedade de materiais, tendo como premissa a promoção do engajamento discente, a articulação entre teoria e prática e o respeito ao ritmo e à autonomia do discente no contexto da EaD. Como exemplos de materiais a serem utilizados no curso de CDIA, destacam-se: textos digitais, livros eletrônicos, videoaulas, podcasts, objetos de aprendizagem interativos, infográficos, simulações, atividades práticas orientadas, entre outros recursos educacionais digitais.

Os AVAs são ambientes essenciais para favorecer o engajamento dos discentes, promovendo o compartilhamento de diferentes linguagens e formatos de materiais didáticos. No curso de CDIA, cada unidade curricular, ofertada no formato a distância, contará com sua própria sala virtual que permitirá o compartilhamento de materiais didáticos e uma comunicação mais eficiente e ágil entre todos os envolvidos nos processos de ensino e aprendizagem. Na UTFPR, os AVAs institucionais são o Moodle e o Google Classroom.

As atividades avaliativas, das unidades curriculares ofertadas integral ou parcial a distância no curso de CDIA, devem estar alinhadas ao planejamento e à organização pedagógica do curso, promovendo uma avaliação contextualizada e diversificada que articule teoria e prática e considere as competências e vivências dos discentes, por meio do uso de múltiplos instrumentos avaliativos. Cada unidade curricular deverá prever, no mínimo, uma atividade avaliativa presencial, com peso majoritário na composição da nota final e o desenvolvimento de habilidades discursivas de análise e síntese. Para complementação da nota da unidade curricular, a critério do docente formador, poderão ocorrer outras atividades avaliativas de forma assíncrona ou síncrona.

O curso de CDIA entende que a adoção de unidades curriculares no formato a distância em cursos de graduação presencial reafirma o compromisso institucional com a inovação pedagógica, a flexibilização curricular responsável e a ampliação das possibilidades formativas, em estrita conformidade com a legislação vigente \cite{UTFPR_EAD_2022, Analise_DNC_2021}. Essa abordagem contribui para a formação integral do discente, preparando-o para atuar de maneira crítica, ética e competente em contextos profissionais e sociais cada vez mais mediados pela TDIC.

Diante desse cenário, a adoção do formato a distância nos cursos vinculados ao Programa Universidades Inovadoras, especialmente nos projetos de curso do $STEM^{IA}$ da UTFPR, amplia as estratégias formativas, fortalece o desenvolvimento de competências tecnológicas e promove maior flexibilidade pedagógica. Tal abordagem não descaracteriza o compromisso do curso de CDIA, com uma formação integral,

Diante desse cenário, a adoção do formato a distância nos cursos vinculados ao Programa Universidades Inovadoras, especialmente nos projetos de curso do $STEM^{IA}$ da UTFPR, amplia as estratégias formativas, fortalece o desenvolvimento de competências tecnológicas e promove maior flexibilidade pedagógica. Tal abordagem não descaracteriza o compromisso do curso de CDIA, com uma formação integral, preservando o rigor acadêmico, mas ao mesmo tempo contemplará a ética, a cidadania, a tecnologia, a sustentabilidade, a cultura e a inovação, assegurando a indissociabilidade entre ensino, pesquisa e extensão.

\section{Concepção Metodológica e Estratégias de Ensino Aprendizagem}
\label{sec:metodologia}

As concepções metodológicas do curso de CDIA, estruturado no modelo $STEM^{IA}$, fundamentam-se na integração entre ciência, tecnologia e prática profissional, valorizando a aprendizagem ativa, investigativa e situada. A formação é guiada por princípios da educação tecnológica, pela articulação contínua entre teoria e prática e pelo protagonismo do estudante como sujeito autônomo em seu processo formativo.

O curso se organiza a partir de práticas pedagógicas que estimulam a autonomia intelectual, a capacidade de tomada de decisão e o desenvolvimento progressivo de competências profissionais. Os estudantes são incentivados a planejar, monitorar e avaliar seus próprios percursos de aprendizagem, de modo a desenvolver independência cognitiva, responsabilidade acadêmica e domínio reflexivo sobre saberes técnicos, científicos e éticos. Essa autonomia, no entanto, é orientada por processos formativos qualificados, nos quais o docente atua como mediador, pesquisador de sua prática e profissional com liberdade pedagógica para selecionar métodos, recursos, sequências didáticas e tecnologias educacionais.

A articulação entre teoria e prática aparece como eixo estruturante do curso. As unidade curricular s de base científica e tecnológica, como matemática, estatística, lógica, programação, sistemas inteligentes e ciência de dados, são ofertadas em ambientes equipados, com laboratórios especializados, plataformas de simulação, ambientes virtuais de desenvolvimento e recursos digitais avançados que permitem experimentar algoritmos, testar modelos, manipular dados e avaliar soluções em diferentes níveis de complexidade. Essa integração fortalece a compreensão dos conceitos e possibilita que o estudante vivencie situações reais desde os primeiros períodos.

Em unidades curriculares aplicadas, atividades como hackathons, oficinas temáticas, desafios de inovação, resolução de problemas reais, desenvolvimento de protótipos e projetos orientados consolidam o conhecimento teórico e fomentam a criatividade, a análise crítica e a responsabilidade técnica. Visitas técnicas, interação com ecossistemas de inovação, projetos integradores e parcerias com empresas e organizações do setor produtivo ampliam o contato com práticas profissionais e com demandas tecnológicas contemporâneas.

O curso incorpora metodologias ativas — Aprendizagem Baseada em Projetos (ABP), Aprendizagem Baseada em Problemas (PBL), sala de aula invertida, projetos colaborativos, prototipagem rápida, uso intensivo de ambientes digitais — que fortalecem o protagonismo discente e a aprendizagem significativa. Tais estratégias permitem que os estudantes articulem conhecimentos de diferentes áreas STEM, desenvolvendo pensamento crítico, criatividade, autonomia e capacidade de resolução de problemas complexos.

A acessibilidade também constitui como orientação central da prática pedagógica. O curso garante condições de aprendizagem inclusivas, com uso de tecnologias assistivas, materiais acessíveis, ambientes adequados e flexibilização de estratégias de ensino, respeitando diferentes ritmos, estilos de aprendizagem e necessidades específicas. A acessibilidade é compreendida não apenas como adequação física ou tecnológica, mas como compromisso ético com a democratização do conhecimento e com a permanência estudantil, em consonância com as políticas institucionais da UTFPR.

A autonomia docente é assegurada como princípio indispensável à qualidade pedagógica. Os professores têm liberdade para definir metodologias, planejar recursos, construir práticas inovadoras e selecionar abordagens didáticas alinhadas ao perfil da turma e aos objetivos de cada unidade curricular . Essa autonomia é acompanhada de responsabilidade acadêmica, compromisso com o PPI \cite{UTFPR_PPI_2019}, com o PDI \cite{UTFPR_PDI_2017, UTFPR_PDI_2023} e com as diretrizes deste PPC, favorecendo um ambiente de ensino dinâmico, crítico e permanentemente atualizado.

A curricularização da extensão fortalece ainda mais a conexão entre teoria, prática e impacto social. Projetos extensionistas em IA aplicados a problemas reais, como análise de dados para políticas públicas, soluções inteligentes para demandas comunitárias, automação de processos em organizações locais ou desenvolvimento de aplicações de acessibilidade, ampliam o compromisso ético e público da formação tecnológica. Esses projetos aproximam o estudante dos Objetivos de Desenvolvimento Sustentável (ODS) e reforçam a compreensão de que a tecnologia deve servir ao bem comum, à inclusão e ao desenvolvimento sustentável \cite{UNICEF_ODS_2030}.

Assim, as estratégias metodológicas do curso buscam formar profissionais críticos, autônomos, criativos e tecnicamente competentes, capazes de compreender, projetar, implementar e avaliar soluções de Inteligência Artificial em diferentes contextos. Ao promover a articulação entre teoria e prática, assegurar condições de acessibilidade, valorizar a autonomia docente e estimular a autonomia discente, este curso consolida uma formação alinhada às demandas contemporâneas da transformação digital e aos valores da educação tecnológica comprometida com a inovação e com a sociedade.

O curso de CDIA adota um Programa de Gestão da Aprendizagem institucionalizado, documentado no Projeto Pedagógico do Curso (PPC), com cronograma definido e efetivamente implementado, em consonância com as normativas acadêmicas da UTFPR. O Programa constitui um instrumento sistemático de planejamento, acompanhamento, avaliação e aprimoramento do processo de ensino-aprendizagem ao longo de todo o percurso formativo, orientando decisões pedagógicas fundamentadas em evidências.

O Programa de Gestão da Aprendizagem conta com ampla e relevante participação do corpo docente do curso, sendo supervisionado, validado e acompanhado pelo Núcleo Docente Estruturante (NDE), com apoio da coordenação e do colegiado do curso. Sua implementação está integrada aos processos de avaliação institucional e aos fluxos regulares de planejamento acadêmico, fortalecendo a gestão coletiva do currículo e a coerência entre o perfil do egresso, as competências previstas no PPC e as práticas pedagógicas adotadas.

No âmbito da aprendizagem discente, o Programa contempla mecanismos sistemáticos de avaliação ao longo do desenvolvimento das unidades curriculares e em relação ao conjunto das competências gerais e específicas do egresso. São utilizados instrumentos quantitativos e qualitativos, incluindo avaliações diagnósticas, formativas e somativas, que permitem o acompanhamento contínuo da aprendizagem, a identificação de lacunas formativas e a análise do desenvolvimento progressivo das competências ao longo do curso.

O Programa de Gestão da Aprendizagem prevê, ainda, a concepção, implementação e avaliação de intervenções pedagógicas, curriculares ou metodológicas destinadas à mitigação das lacunas identificadas. Para isso, são utilizados instrumentos periódicos de aferição que comparam a aprendizagem prevista em cada etapa do percurso formativo com a aprendizagem efetivamente alcançada. Os resultados dessas análises subsidiam a avaliação continuada das intervenções realizadas e orientam o aprimoramento permanente do currículo e do PPC, assegurando a qualidade acadêmica e a aderência do curso às diretrizes institucionais e às demandas contemporâneas da formação tecnológica.

\section{Procedimentos de Avaliação da Aprendizagem e Recuperação de Aprendizagem}

No curso de CDIA do campus Londrina, estruturado no modelo $STEM^{IA}$ da UTFPR, o processo de avaliação da aprendizagem orienta-se pelos princípios da formação tecnológica, da indissociabilidade entre teoria e prática, e da promoção de autonomia intelectual do estudante. A avaliação é compreendida como parte integrante do processo formativo e tem caráter contínuo, diagnóstico e formativo, buscando identificar avanços, dificuldades e necessidades de reorganização das estratégias pedagógicas.

Os docentes utilizam um conjunto diversificado de instrumentos avaliativos, selecionados em consonância com as competências e habilidades previstas no PPC e com a natureza específica de cada unidade curricular. Entre as estratégias adotadas, destacam-se: provas escritas, atividades práticas em laboratório, desenvolvimento e apresentação de projetos, estudos de caso, desafios de programação, seminários, listas de exercícios, relatórios técnicos, participação em atividades extensionistas e avaliações processuais desenvolvidas ao longo do período letivo. A definição dos critérios, pesos e instrumentos é de responsabilidade do docente, respeitando sua autonomia pedagógica, e é apresentada aos estudantes nos planos de ensino, garantindo transparência e previsibilidade ao processo avaliativo.

Considerando a característica interunidade curricular e aplicada do modelo $STEM^{IA}$, muitas avaliações incluem resolução de problemas reais, integração entre diferentes áreas de conhecimento, uso de ambientes computacionais e elaboração de soluções baseadas em algoritmos de IA, análise de dados ou prototipagem digital. Esses instrumentos favorecem a articulação entre teoria e prática, estimulam o pensamento crítico, a criatividade e a capacidade de tomada de decisão fundamentada — competências essenciais ao perfil do egresso.

A recuperação da aprendizagem é garantida como direito do estudante e ocorre de forma contínua, acompanhando o desenvolvimento das atividades de cada unidade curricular. Os docentes podem oferecer diferentes oportunidades, como atividades de reforço, tarefas complementares, reavaliações, projetos adicionais, estudos orientados ou atendimentos individuais, sempre com o objetivo de apoiar o estudante na superação de dificuldades e no alcance dos resultados de aprendizagem esperados. Essas ações devem ser planejadas no decorrer do período letivo, evitando processos concentrados ou meramente compensatórios, reafirmando a natureza pedagógica da avaliação.

No âmbito institucional, a UTFPR mantém processos regulares de avaliação de unidade curricular s e cursos, que envolvem estudantes, docentes e coordenação. O Núcleo Docente Estruturante (NDE), a coordenação do curso e as instâncias institucionais de apoio psicopedagógico dispõem de ferramentas analíticas e relatórios acadêmicos que permitem acompanhar a progressão do desempenho dos estudantes, identificar lacunas de aprendizagem e subsidiar decisões pedagógicas e curriculares.

Os resultados dessas avaliações são utilizados tanto para o aprimoramento contínuo dos planos de ensino quanto para subsidiar ações de formação continuada do corpo docente, contribuindo para o aperfeiçoamento das práticas avaliativas, metodológicas e pedagógicas adotadas no curso. A participação da comunidade acadêmica nessas avaliações fortalece a gestão democrática, reforça a transparência e contribui para a consolidação de uma cultura de qualidade na formação tecnológica.

Assim, o processo de avaliação no curso $STEM^{IA}$ articula rigor acadêmico, acompanhamento pedagógico e oportunidades de recuperação da aprendizagem, assegurando a formação de profissionais críticos, éticos, criativos e tecnicamente qualificados para atuar nos diversos contextos da Inteligência Artificial e da transformação digital.

\section{Processo de Aprendizagem Mediada por Tecnologia e Tecnologias Utilizadas}

No curso de CDIA do \textit{campus} Londrina, os processos de aprendizagem mediados por tecnologias estruturam a proposta pedagógica e acompanham a centralidade das TDIC e da Inteligência Artificial na formação contemporânea. As tecnologias são tratadas simultaneamente como conteúdo e como instrumento didático, favorecendo a articulação entre teoria e prática, a resolução de problemas complexos, o desenvolvimento de projetos, a autonomia do estudante e o fortalecimento do pensamento computacional.

A utilização de ambientes digitais, como Moodle institucional, \textit{Google Workspace}, laboratórios virtuais, simuladores, ferramentas de análise de dados e plataformas de desenvolvimento, ocorre de forma transversal nas unidades curriculares e nas atividades de ensino, pesquisa e extensão. Esses recursos ampliam a experimentação, diversificam as formas de interação e garantem acessibilidade digital, continuidade das atividades acadêmicas e integração entre momentos síncronos e assíncronos.

As estratégias pedagógicas apoiadas por tecnologias incluem metodologias ativas, estudos de caso digitais, projetos colaborativos, experimentação em ambientes simulados e acompanhamento contínuo da aprendizagem por meio de recursos digitais. Tais práticas promovem letramento digital, autonomia intelectual e preparo para atuar em ecossistemas profissionais intensivos em dados, IA, automação e inovação tecnológica.

Assim, a mediação tecnológica assegura coerência entre o perfil do egresso, as demandas da sociedade digital e o compromisso institucional da UTFPR com uma formação sólida, atualizada e orientada pelos avanços da ciência, tecnologia e inteligência artificial.

\section{Tecnologias Digitais no Âmbito do Curso}

O curso de CDIA da UTFPR – campus Londrina fundamenta sua organização didático-pedagógica no uso intensivo e crítico de tecnologias digitais, coerente com a natureza formativa da área e com as demandas contemporâneas do mundo do trabalho. As atividades acadêmicas são apoiadas pelo Ambiente Virtual de Aprendizagem (AVA) institucional, baseado na plataforma Moodle, bem como por suítes colaborativas digitais, que viabilizam comunicação síncrona e assíncrona, compartilhamento de materiais, desenvolvimento de projetos em equipe e acompanhamento contínuo do desempenho discente.

A infraestrutura do curso contempla laboratórios equipados com computadores de alto desempenho, recursos de processamento para análise de grandes volumes de dados e execução de modelos de Inteligência Artificial, além de softwares e bibliotecas amplamente utilizados no contexto profissional e científico. Essa estrutura possibilita a realização de atividades práticas, projetos integradores, experimentação com dados reais e desenvolvimento de soluções computacionais alinhadas às práticas do setor produtivo.

O uso das tecnologias digitais é orientado por princípios de acessibilidade, inclusão e ética no tratamento de dados. O curso conta com suporte institucional para a promoção do letramento digital de docentes e técnicos administrativos, bem como com ações de formação continuada voltadas ao uso pedagógico de ferramentas digitais e sistemas baseados em Inteligência Artificial. Recursos de acessibilidade digital e tecnologias assistivas são disponibilizados conforme as necessidades dos estudantes, assegurando equidade no processo formativo.

Os materiais didáticos priorizam formatos digitais, interativos e atualizados, sendo avaliados pelo Núcleo Docente Estruturante quanto à sua pertinência, atualidade e alinhamento ao perfil do egresso. O AVA é utilizado como espaço de organização das unidades curriculares, disponibilização de conteúdos, aplicação de atividades avaliativas e devolutivas formativas, fortalecendo o protagonismo discente e a transparência acadêmica.

Além disso, considerando a especificidade do curso, ferramentas de programação colaborativa, ambientes de versionamento de código, plataformas de computação em nuvem e recursos de Inteligência Artificial são incorporados às práticas pedagógicas, sempre sob orientação docente, promovendo experiências formativas alinhadas às exigências tecnológicas contemporâneas.

A efetividade das tecnologias adotadas é monitorada por meio de processos institucionais de avaliação, incluindo a atuação da Comissão Própria de Avaliação (CPA) e consultas periódicas à comunidade acadêmica, assegurando atualização contínua, adequação pedagógica e aderência às transformações tecnológicas da área.

\section{Processo de Gestão do Curso}

O processo de gestão do curso, descrito neste PPC, estrutura-se de forma sistemática, participativa e orientada por evidências, com responsabilidades definidas entre Coordenação de Curso, Núcleo Docente Estruturante (NDE) e Colegiado, articulando-se à política de avaliação institucional da UTFPR conduzida pela Comissão Própria de Avaliação (CPA). A gestão utiliza múltiplas fontes de informação para monitoramento e tomada de decisão, incluindo: relatórios e indicadores da autoavaliação institucional (CPA), resultados de avaliações externas (MEC/INEP), avaliação de componentes curriculares e do docente pelo discente, indicadores acadêmicos (ingresso, evasão, retenção, integralização, desempenho), acompanhamento de egressos e manifestações oriundas de canais institucionais de comunicação e transparência.

A autoavaliação do curso possui periodicidade definida e integra um ciclo contínuo de análise, planejamento e replanejamento. Com base nas evidências coletadas, a Coordenação e o NDE realizam análises periódicas para identificar oportunidades de melhoria e, quando necessário, propor ajustes na estrutura curricular, metodologias e processos avaliativos, assegurando o aprimoramento da eficácia do currículo no desenvolvimento do perfil do egresso e das competências previstas. Essas análises resultam em plano de ação formal, com definição de responsáveis, cronograma, metas e indicadores de acompanhamento, bem como avaliação subsequente da eficácia das ações implementadas, garantindo melhoria contínua e rastreabilidade das decisões.

Os resultados das avaliações internas e externas são disponibilizados à comunidade acadêmica por meio de tecnologias digitais institucionais, assegurando acesso, transparência e participação, considerando todos os locais de oferta e ambientes formativos vinculados ao curso (conforme o caso). Sempre que pertinente, informações consolidadas também são divulgadas ao público externo, reforçando a prestação de contas e a visibilidade do curso.

O acompanhamento e monitoramento dos resultados avaliativos é realizado de forma contínua pela gestão do curso, em articulação com as instâncias institucionais, utilizando múltiplas evidências para identificar melhorias acadêmicas, administrativas e pedagógicas. Esse processo inclui acompanhamento multidimensional do desempenho docente e, em consonância com as políticas institucionais, subsidia o planejamento e a implementação de ações de capacitação e formação continuada (pedagógica, tecnológica e metodológica), retroalimentando o aprimoramento das práticas docentes e dos planos de ensino.

Dessa forma, o processo de gestão do curso assegura governança acadêmica, transparência, avaliação sistemática e melhoria contínua, fortalecendo a qualidade da formação e a aderência às diretrizes institucionais da UTFPR e às exigências regulatórias do MEC/INEP.

\section{Sistema de Acolhimento e Apoio aos Estudantes}

A UTFPR dispõe de uma estrutura institucional voltada ao acolhimento, permanência e sucesso acadêmico dos estudantes, articulada ao fortalecimento do processo de ensino-aprendizagem e à promoção da inclusão e da equidade. No âmbito dos campi, essa estrutura é coordenada pelo Departamento de Educação (DEPED), vinculado à Diretoria de Graduação e Educação Profissional (DIRGRAD), responsável por ações pedagógicas e de apoio psicossocial.

O DEPED é constituído pelo Núcleo de Ensino (NUENS), voltado à gestão pedagógica e ao assessoramento docente, e pelo Núcleo de Acompanhamento Psicopedagógico e Assistência Estudantil (NUAPE), responsável pelo atendimento coletivo e individualizado aos discentes. O NUAPE desenvolve ações de acompanhamento psicopedagógico, assistência estudantil e apoio à permanência, atuando de forma integrada com as coordenações de curso e com os demais setores institucionais.

Em consonância com a Lei nº 13.146, de 06 de julho de 2015 (Lei Brasileira de Inclusão da Pessoa com Deficiência) \cite{Brasil_2015_Deficiencia}, e com o Decreto nº 7.611, de 17 de novembro de 2011 \cite{Brasil_2011_Deficiencia}, que dispõe sobre a educação especial e o atendimento educacional especializado, a UTFPR assegura condições de acesso, participação e aprendizagem aos estudantes com deficiência, transtornos globais do desenvolvimento, altas habilidades/superdotação e demais necessidades educacionais específicas.

No campus Londrina, o NUAPE integra o Núcleo de Acessibilidade e Inclusão (NAI) e o setor de Assistência à Saúde. O NAI promove ações voltadas à acessibilidade pedagógica, arquitetônica, comunicacional e digital, apoiando docentes na adequação de estratégias metodológicas e avaliativas, quando necessário. Atua, ainda, na promoção de uma cultura institucional orientada à convivência com a diversidade e ao respeito às diferenças, articulando recursos humanos, materiais e tecnológicos adequados às demandas identificadas.

O setor de Assistência à Saúde oferece atendimento de enfermagem e acompanhamento psicológico educacional, contribuindo para o bem-estar físico e emocional dos estudantes. A equipe multiprofissional do campus é composta por psicólogas, pedagoga, assistentes sociais e profissional de apoio em enfermagem, que atuam tanto de forma preventiva quanto interventiva, conforme as necessidades apresentadas.

O NUAPE é responsável por operacionalizar os programas de Auxílio Estudantil e Bolsa Permanência do MEC, que incluem as seguintes modalidades:

\begin{itemize}
    \item Auxílio estudantil: é um programa institucional que tem por finalidade apoiar o discente para o desenvolvimento acadêmico e sua permanência na instituição, buscando reduzir os índices de evasão, decorrentes de dificuldades de ordem socioeconômicas. O Auxílio Estudantil é dividido em quatro modalidades: Auxílio Básico, Auxílio Alimentação, Auxílio Moradia e Auxílio Instalação.
    \item Bolsa Permanência do MEC: é uma ação do governo federal de concessão de auxílio financeiro aos estudantes em situação de vulnerabilidade socioeconômica e para estudantes indígenas e quilombolas. Tem por finalidade minimizar as desigualdades sociais e contribuir para a permanência e a diplomação dos estudantes de graduação.
\end{itemize}

O apoio ao estudante no curso de CDIA também se concretiza por meio do acompanhamento sistemático da coordenação de curso, do Núcleo Docente Estruturante (NDE) e do colegiado, que analisam indicadores acadêmicos, taxas de retenção e evasão, bem como demandas específicas identificadas ao longo do percurso formativo. Sempre que necessário, são encaminhadas ações conjuntas com o NUAPE e demais setores institucionais.

Adicionalmente, o curso promove ações de acolhimento aos ingressantes, orientação acadêmica sobre organização dos estudos, uso de tecnologias digitais e compreensão do projeto pedagógico, favorecendo a adaptação ao ensino superior e à dinâmica própria da área de Ciência de Dados e Inteligência Artificial.

Dessa forma, a UTFPR campus Londrina reafirma seu compromisso com a formação integral do estudante, assegurando condições institucionais para acesso, permanência qualificada, desenvolvimento acadêmico e inserção profissional, em um ambiente universitário pautado pela inclusão, diversidade e respeito às diferenças.

\section{Ações de Apoio à Inserção Profissional}

A UTFPR desenvolve ações estruturadas de apoio à inserção profissional dos estudantes, que são planejadas, acompanhadas e avaliadas com a participação do Núcleo Docente Estruturante (NDE) e da Coordenação de Curso, assegurando coerência entre formação acadêmica, perfil do egresso e demandas do mundo do trabalho.

A instituição mantém e amplia parcerias formais com empresas, órgãos públicos, instituições de pesquisa, parques tecnológicos, incubadoras e arranjos produtivos locais, como descrito na seção Integração do curso com o mundo do trabalho (Seção \ref{sec:integracaoMundoDoTrabalho}), viabilizando estágios, projetos aplicados, visitas técnicas, desafios de inovação e competições científicas. No âmbito dos cursos $STEM^{IA}$, essas parcerias assumem papel estratégico, promovendo vivências em ambientes de ciência de dados, automação e Inteligência Artificial, e subsidiando a formulação de atividades práticas, projetos integradores e ações extensionistas.

As ações de apoio à inserção profissional incluem sistemas institucionais de cadastro e monitoramento de oportunidades, estágios e trajetórias de egressos, bem como a divulgação regular dessas oportunidades por canais institucionais ágeis, transparentes e amplamente conhecidos pela comunidade acadêmica. Esses dados subsidiam a celebração de convênios, ações de mobilidade estudantil e o aprimoramento contínuo do curso.

As práticas profissionais supervisionadas são desenvolvidas em ambientes reais do exercício profissional, com coordenação formalmente designada pela instituição e acompanhamento por docentes do curso e profissionais das organizações parceiras. Essas práticas proporcionam ensino contextualizado, promovem a interunidade curricular ridade e a integração entre ensino, pesquisa e extensão, preparando os estudantes para desafios interunidade curricular res e demandas contemporâneas da sociedade.

Os resultados das práticas e projetos são, quando pertinente, publicizados para a comunidade acadêmica e externa, fortalecendo a transparência, a integração com o mundo do trabalho e a formação profissional qualificada dos estudantes dos cursos $STEM^{IA}$.

A UTFPR desenvolve ações estruturadas de apoio à inserção profissional dos estudantes, articulando formação acadêmica, práticas profissionais e aproximação com o mundo do trabalho. Para fortalecer esse processo, a instituição mantém e amplia parcerias com empresas, órgãos públicos, instituições de pesquisa, parques tecnológicos, incubadoras e arranjos produtivos locais, possibilitando oportunidades de estágios, projetos aplicados, programas de residência tecnológica, visitas técnicas e participação em desafios e competições científicas. No âmbito dos cursos $STEM^{IA}$, essas parcerias assumem papel estratégico, promovendo a vivência em ambientes de inovação, ciência de dados, automação e Inteligência Artificial, favorecendo a transição qualificada para o mercado profissional.

Além das articulações externas, a UTFPR oferece ações internas de orientação e acompanhamento para inserção profissional, incluindo apoio ao desenvolvimento de competências socioemocionais, atividades de capacitação, cursos livres, palestras, feiras de carreiras, mentorias, oficinas de elaboração de currículo e preparação para processos seletivos. A atuação integrada da Coordenação de Curso, do Núcleo Docente Estruturante e dos setores institucionais responsáveis pela carreira e empreendedorismo contribui para formar profissionais capazes de dialogar com as demandas tecnológicas emergentes, ampliando as condições de empregabilidade e protagonismo dos egressos do curso de CDIA.

A integração do CDIA com a sociedade constitui princípio estruturante da proposta formativa e está explicitamente contemplada neste Projeto Pedagógico do Curso (PPC). O curso integra, de forma transversal e interunidade curricular, os temas de sustentabilidade, diversidade e responsabilidade social, articulando-os às competências técnicas e profissionais próprias da formação em áreas STEM com ênfase em Inteligência Artificial. Esses temas são abordados ao longo do percurso formativo, tanto nas unidades curriculares regulares quanto nas ações de extensão curricularizadas, contribuindo para a formação cidadã, ética e socialmente responsável do egresso.

A abordagem desses temas ocorre de maneira contextualizada e aplicada, considerando os impactos sociais, ambientais, éticos e econômicos das tecnologias digitais, da ciência de dados e da inteligência artificial por meio de disciplinas do Ciclo de Humanidades e também pelas disciplinas Projeto Extensionista e Oficinas de Integração 1 e 2. O curso promove reflexões e práticas que problematizam o uso responsável da tecnologia, a redução de desigualdades, a inclusão digital, a acessibilidade, a proteção de dados, a equidade e a sustentabilidade, assegurando coerência entre formação técnica, compromisso social e valores institucionais da UTFPR.

A integração com a sociedade é operacionalizada por meio de ações, programas, projetos, cursos, oficinas, eventos de extensão e prestação de serviços tecnológicos, desenvolvidos de forma regular e articulados às demandas do território em que o curso está inserido. Com o apoio da Diretoria de Relações Empresariais e Comunitárias (DIREC) o curso busca alinhar suas disciplinas extensionistas, bem como atividades como estágios e TCCs às demandas reais da comunidade local. Essas ações envolvem a participação ativa de docentes, estudantes e parceiros externos, incluindo organizações públicas, privadas, do terceiro setor, ambientes profissionais, parques tecnológicos e arranjos socioprodutivos locais e regionais.

As ações extensionistas e de interação social são avaliadas sistematicamente pelo Núcleo Docente Estruturante (NDE) e pelo Colegiado do Curso, assegurando alinhamento com o perfil do egresso, com os objetivos formativos e com as competências previstas no PPC. Tais ações promovem o reconhecimento da função social do profissional formado, qualificam a pesquisa aplicada (quando pertinente) e contribuem para o desenvolvimento social, tecnológico e sustentável das comunidades atendidas.
Dessa forma, o curso de CDIA reafirma seu compromisso com uma formação tecnológica integrada à realidade social, fortalecendo o diálogo entre universidade e sociedade e consolidando o papel do profissional STEMIA como agente de inovação responsável, transformação social e desenvolvimento sustentável

\section{Perfil do Egresso}

O perfil do egresso do curso de CDIA está alinhado aos referenciais para formação de Cursos de Bacharelado em Ciência de Dados \cite{SBC_REF_CD_2023} e de formação de Cursos de Bacharelado em Inteligência Artificial \cite{SBC_REF_IA_2024}.

Na linha de Ciência de Dados, o curso visa formar profissionais capazes de ``pensar com dados''. Esses profissionais devem possuir competência teórica, técnica e metodológica, bem como experiência prática para lidar com as mais variadas situações e domínios de aplicação. Espera-se que os egressos sejam capazes de:

\begin{itemize}
    \item Entender, formular e refinar as questões apropriadas no processo de análise de dados;
    \item Obter, modelar e explorar os dados relacionados;
    \item Processar os dados e realizar as análises necessárias;
    \item Obter e comunicar o conhecimento relevante da área;
    \item Apoiar o desenvolvimento e implantação de soluções com base nos resultados atingidos; e,
    \item Entender e atender aspectos éticos e sociais relacionados à sua atuação.
\end{itemize}

Na linha de Inteligência Artificial espera-se que o egresso:

\begin{itemize}
    \item Tenha sólida formação em computação, matemática e estatística que os capacitem a construir soluções computacionais para problemas complexos;
    \item Conheça os principais paradigmas de sistemas de IA e os processos envolvidos na sua construção e análise;
    \item Domine o uso de técnicas de aquisição, tratamento, mineração e visualização de dados e de algoritmos de aprendizado de máquina
    \item Seja capaz de criar soluções, individualmente ou em equipe, para problemas complexos e reconheça o caráter fundamental da inovação e da criatividade e compreender as perspectivas de negócios e oportunidades comerciais relevantes na área de Inteligência Artificial;
    \item Seja capaz de agir de forma consciente na construção e no uso de sistemas de IA, compreendendo seus impactos diretos e indiretos sobre a sociedade e os indivíduos no que tange a questões de privacidade, transparência, equanimidade e preconceito implícito em sistemas de software e bases de dados, entre outras.
\end{itemize}

Em síntese, o conjunto de competências e habilidades descritas consolida um perfil de egresso tecnicamente robusto, metodologicamente consistente e socialmente responsável. A formação proposta integra fundamentos científicos, domínio tecnológico e capacidade analítica, articulando Ciência de Dados e Inteligência Artificial de maneira complementar e estratégica.

O egresso estará apto a atuar em diferentes setores econômicos e sociais, desenvolvendo, avaliando e implementando soluções baseadas em dados e sistemas inteligentes, com visão crítica, postura ética e compromisso com a inovação. Dessa forma, o curso contribui para a formação de profissionais preparados para enfrentar desafios complexos, promover transformação digital e gerar impacto positivo no desenvolvimento regional e nacional.

O curso de CDIA se baseia nas competências elencadas nos referenciais para formação de Cursos de Bacharelado em Ciência de Dados \cite{SBC_REF_CD_2023} e de formação de Cursos de Bacharelado em Inteligência Artificial \cite{SBC_REF_IA_2024}.

Considera-se competência como capacidade de mobilizar, integrar e aplicar saberes teóricos, práticos, socioemocionais e contextuais para resolver problemas reais em ambientes complexos e em transformação, com progressão de complexidade ao longo do percurso formativo, conforme a concepção adotada no PPC.

Na linha de Ciência de Dados, os referenciais incluem 8 competências do egresso, a saber: (i) fundamentos de matemática, estatística e computação para ciência de dados, (ii) modelagem estatística e resolução de problemas, (iii) desenvolvimento de sistemas; (iv) engenharia e exploração de dados, (v) big data, (vi) mineração de dados e aprendizado de máquina, (vii) aprendizado contínuo e autônomo e (viii) ciência, tecnologia e inovação.

Por outro lado, na linha de Inteligência Artificial, os referenciais recomendam que o egresso desenvolva 7 competências, sendo elas: (i) fundamentos de matemática, estatística e ciência da computação, (ii) desenvolvimento e gestão de sistemas de IA, (iii) raciocínio e representação do conhecimento, (iv) ciência de dados, (v) aprendizado de máquina, (vi) percepção e atuação: visão computacional, processamento de linguagem natural e robótica e (vii) aperfeiçoamento pessoal e profissional.

Considerando que há sobreposições e complementariedades entre as competências de ambos referenciais, e também considerando as características regionais e do corpo docente do curso, foram formuladas 8 competências para o curso de CDIA, as quais são listadas no Quadro \ref{quadro:2.13}.

\begin{quadro}[htb]
    \centering
    \caption{Relação de competências desenvolvidas ao longo do curso}
    \label{quadro:2.13}
    \begin{tabular}{|c|p{4cm}|p{10cm}|}
        \hline
        \textbf{Código} & \textbf{Competência} & \textbf{Descrição} \\ \hline
        C1 & Fundamentos matemáticos, estatísticos e computacionais & Compreender e aplicar fundamentos de matemática, estatística, algoritmos e ciência da computação na modelagem e resolução de problemas em Ciência de Dados e Inteligência Artificial, considerando limites teóricos e computacionais. \\ \hline
        C2 & Modelagem, análise e solução de problemas complexos & Analisar problemas complexos e desenvolver modelos analíticos, estatísticos e computacionais para apoiar a tomada de decisão, integrando dados, conhecimento e técnicas de Inteligência Artificial. \\ \hline
        C3 & Engenharia de dados e computação em larga escala & Projetar, implementar e gerenciar soluções de armazenamento, integração, governança e processamento de dados, incluindo Big Data, assegurando qualidade, segurança, privacidade e conformidade legal. \\ \hline
        C4 & Aprendizado de máquina e métodos de Inteligência Artificial & Projetar, implementar, avaliar e aplicar técnicas de aprendizado de máquina, aprendizado profundo e raciocínio simbólico, considerando desempenho, interpretabilidade, vieses e aplicações reais. \\ \hline
        C5 & Sistemas inteligentes e aplicações & Projetar, desenvolver e integrar sistemas inteligentes completos, incluindo percepção, linguagem, interação humano-IA e atuação em ambientes físicos ou digitais, assegurando qualidade, escalabilidade e usabilidade. \\ \hline
        C6 & Ética, impacto social e responsabilidade profissional & Analisar criticamente os impactos sociais, econômicos, ambientais e culturais da Ciência de Dados e da Inteligência Artificial, atuando de forma ética, responsável e alinhada aos marcos legais e regulatórios. \\ \hline
        C7 & Inovação, empreendedorismo e gestão & Atuar de forma inovadora e empreendedora no desenvolvimento de soluções baseadas em dados e Inteligência Artificial, aplicando princípios de gestão, liderança, planejamento estratégico e viabilidade técnica e econômica. \\ \hline
        C8 & Comunicação, colaboração e aprendizado contínuo & Comunicar ideias e resultados de forma clara para públicos técnicos e não técnicos, atuar em equipes multidisciplinares, gerenciar projetos e manter aprendizado contínuo frente à evolução tecnológica. \\ \hline
    \end{tabular}
    \fonte{Autoria prória (2026)}
\end{quadro}

\subsection{Reconhecimento de saberes, competências e aproveitamento de estudo}

O curso de CDIA adota diretrizes institucionais alinhadas ao Plano de Desenvolvimento Institucional da UTFPR, assegurando mecanismos formais para o reconhecimento de saberes, competências e o aproveitamento de estudos prévios, conforme legislação educacional vigente. Esses mecanismos são estruturados a partir de evidências objetivas de aprendizagem, permitindo a validação de conhecimentos adquiridos em contextos acadêmicos, profissionais e formativos diversos.

A organização curricular prevê, em seu itinerário formativo, procedimentos sistematizados para o reconhecimento de competências provenientes de experiências de trabalho, cursos técnicos, cursos de nível tecnológico e certificações profissionais. A avaliação dessas competências ocorre mediante instrumentos que garantem rastreabilidade e confiabilidade, tais como: análise documental, demonstrações práticas, portfólios, entrevistas técnicas, provas de desempenho e correlação direta com os resultados de aprendizagem previstos nas unidades curriculares do curso.

O curso também contempla dispositivos de certificação profissional intermediária, alinhados à política institucional de formação por competências e às diretrizes legais aplicáveis. Essas certificações seguem padrões de verificação baseados em evidências, assegurando que o estudante demonstre domínio efetivo dos conhecimentos e habilidades correspondentes aos marcos formativos estabelecidos no PPC. Tais instrumentos reforçam a empregabilidade e promovem a articulação com setores produtivos intensivos em tecnologia e inovação.

Quanto ao aproveitamento de estudos, o curso estabelece critérios objetivos baseados na correspondência entre ementas, cargas horárias, competências e níveis de complexidade cognitiva, conforme matriz curricular e normas institucionais. A equivalência entre componentes curriculares é avaliada por comissões específicas, garantindo rigor técnico, coerência acadêmica e aderência ao perfil profissional do egresso.

Assim, o Bacharelado em CDIA estrutura-se com base em práticas acadêmicas verificáveis e orientadas por evidências, assegurando processos robustos de reconhecimento de saberes e competências, promovendo itinerários formativos flexíveis e fortalecendo a qualidade da formação tecnológica ofertada pela UTFPR.

\section{Acompanhamento dos Egressos no Mercado de Trabalho}

A UTFPR realiza o acompanhamento sistemático dos egressos com o objetivo de monitorar sua inserção e trajetória no mercado de trabalho, bem como avaliar a aderência entre a formação ofertada e as demandas profissionais dos setores produtivos. Esse processo envolve a aplicação periódica de pesquisas institucionais, o diálogo com empresas parceiras, a análise de indicadores de empregabilidade e a participação dos egressos em atividades acadêmicas, como palestras e ações de atualização profissional.

Esse acompanhamento é realizado pela atuação da Diretoria de Relações Empresariais e Comunitárias (DIREC), que promove a interação entre a Universidade, o setor produtivo e a comunidade, viabilizando parcerias institucionais e empresariais, convênios, eventos, intercâmbios, além do agenciamento de estágios e oportunidades de emprego, contribuindo diretamente para a inserção profissional dos egressos.

Complementarmente, o Departamento de Apoio e Projetos Tecnológicos (DEPET) atua no desenvolvimento de projetos tecnológicos alinhados às demandas do setor produtivo regional, na promoção de ações de inovação, transferência de tecnologia e apoio a parcerias tecnológicas, possibilitando a aproximação dos estudantes e egressos com ambientes reais de inovação, empreendedorismo e desenvolvimento tecnológico.

No contexto do curso de CDIA, esse acompanhamento assume relevância estratégica, ao permitir a identificação de tendências nas áreas de Ciência de Dados, Inteligência Artificial e tecnologias correlatas, retroalimentando o Projeto Pedagógico do Curso e subsidiando ajustes curriculares que fortalecem a formação e a inovação dos futuros profissionais no mercado de trabalho.

\section{Desenvolvimento da Internacionalização}

A internacionalização constitui um dos eixos estratégicos da Universidade Tecnológica Federal do Paraná, orientando-se pelo compromisso de ampliar a formação acadêmica, científica e cultural dos estudantes por meio da integração com instituições de ensino superior do exterior. Coordenada de forma articulada pela Pró-Reitoria de Graduação e Educação Profissional (PROGRAD) e pela Pró-Reitoria de Relações Empresariais e Comunitárias (PROREC), por meio da Diretoria de Relações Interinstitucionais (DIRINTER), a política institucional de internacionalização busca promover experiências formativas que fortaleçam a inserção global da UTFPR e ampliem as oportunidades de mobilidade estudantil, cooperação acadêmica e produção compartilhada de conhecimento.

A internacionalização, nesse contexto, é compreendida como uma dimensão formativa que amplia repertórios, fortalece competências interculturais e impulsiona a qualificação profissional e científica dos estudantes da UTFPR. Ao promover a circulação de pessoas, saberes e práticas, a universidade reafirma seu compromisso com uma formação alinhada às demandas contemporâneas, capaz de integrar perspectivas globais sem perder de vista as necessidades e os desafios do desenvolvimento regional e nacional.

\subsection{Ações de internacionalização}

Entre as ações estruturantes, destaca-se o Programa de Mobilidade Estudantil Internacional (MEI), que possibilita o afastamento temporário de estudantes de graduação para cursar unidade curricular em universidades estrangeiras parceiras, retornando posteriormente à UTFPR para conclusão do curso. A mobilidade é viabilizada por acordos de cooperação que garantem isenção das taxas estritamente acadêmicas nas instituições receptoras, assegurando ao estudante a oportunidade de vivenciar outras culturas acadêmicas, metodologias de ensino e perspectivas de atuação profissional.

O processo de seleção é conduzido com base em princípios de transparência, mérito acadêmico e responsabilidade institucional, considerando requisitos como desempenho acadêmico, proficiência linguística e aderência da proposta de estudos às áreas de formação do estudante. A UTFPR também estabelece procedimentos rigorosos para elaboração, avaliação e validação do plano de estudos, assegurando que a experiência internacional contribua efetivamente para o desenvolvimento acadêmico do aluno.

A matriz curricular deste PPC foi desenvolvida em consonância com o PDI de forma a manter e ampliar a rede de colaboração internacional. Atualmente, a UTFPR possui  convênios estabelecidos com mais de 110 instituições de 24 países (por exemplo, Portugal, 22; França, 11; Alemanha, 9), destacando-se o intercâmbio de estudantes com a França e Portugal (no âmbito da dupla diplomação).

A participação nestes programas de mobilidade internacional proporciona ao estudante o enriquecimento técnico-científico, quer por completar parte do seu curso em outra instituição, quer por realizar estágio curricular em indústrias/empresas estrangeiras, além de  ampliar sua autonomia e fortalecer sua formação cultural e humanística.

Além dessas oportunidades internacionais, é possível ainda incluir a "Internacionalização em Casa" com atividades desenvolvidas em colaboração com instituições estrangeiras por meio de unidades curriculares à distância ou estágio internacional, bem como unidades curriculares a serem ministradas em outros idiomas dentro da própria UTFPR.
